\documentclass[aspectratio=169,10pt]{beamer}
\usetheme{mekatronik}
\setday{1}{Electrical Machines and Actuators}
\title[Electrical Machines and Actuators]{Electrical Machines and Actuators}
\subtitle{Principles, machine structures, motor models, transmissions and sizing\\ + introduction to dynamic system modelling}
\author{Mechatronics Training Programme}
\institute{Theoretical Course Series}
\date{Day 1 of 5}

% small helpers used only in this deck
\newcommand{\figph}[2]{\includegraphics[#1]{../figures/#2}}
\tikzset{
  mkblk/.style={draw=mkBlue,fill=mkPale,rounded corners=1.5pt,minimum height=7mm,minimum width=12mm,font=\scriptsize,align=center,line width=.6pt},
  mksum/.style={draw=mkBlue,circle,fill=white,minimum size=4.5mm,inner sep=0pt,line width=.6pt},
  mksig/.style={-{Stealth[length=2mm]},mkNavy,line width=.6pt},
  mklab/.style={font=\scriptsize,text=mkNavy},
}

\begin{document}

% =====================================================================
\begin{frame}[plain]\titlepage\end{frame}

% =====================================================================
\begin{frame}{Agenda and learning objectives}
\begin{columns}[T]
\begin{column}{.52\textwidth}
\small
\renewcommand{\arraystretch}{1.15}
\begin{tabular}{@{}l p{4.6cm} r@{}}
\toprule
\textbf{Code} & \textbf{Session} & \textbf{Time}\\
\midrule
1.1 & Fundamental principles and classification & 0.50\,h\\
1.2 & Structure of electrical machines & 1.00\,h\\
1.3 & Electrical and mechanical models & 1.00\,h\\
1.4 & Actuators and motion transmission & 0.75\,h\\
1.5 & Characteristics and selection criteria & 0.75\,h\\
2.1 & Modelling of dynamic systems (part 1) & 1.00\,h\\
\midrule
 & \textbf{Net teaching time} & \textbf{5.00\,h}\\
\bottomrule
\end{tabular}\\[1mm]
{\scriptsize\color{mkGray} Breaks ($\approx$10 min) roughly every hour are additional.}
\end{column}
\begin{column}{.45\textwidth}
\begin{defbox}[After today you can \dots]
\footnotesize
\begin{itemize}\setlength\itemsep{1pt}
\item explain how force/torque and voltage arise from $\bm{B}$, $i$ and motion
\item distinguish DC, BLDC, PMSM, induction and stepper machines
\item write and use the DC-motor model, block diagram and transfer function
\item reflect loads through gears and screws; choose a gear ratio
\item size a servo motor from a move profile (peak, RMS, thermal)
\item build lumped models of mechanical, electrical and coupled systems
\end{itemize}
\end{defbox}
\end{column}
\end{columns}
\end{frame}

% =====================================================================
\section{1.1 Fundamental principles and classification}
% =====================================================================

\begin{frame}{Electromechanical energy conversion in a servo axis}
\centering
\begin{tikzpicture}[node distance=5mm and 6mm]
\node[mkblk,minimum width=15mm] (grid) {Supply\\ AC grid / DC};
\node[mkblk,right=of grid,minimum width=15mm] (drv) {Power\\ converter\\ (drive)};
\node[mkblk,right=of drv,fill=mkCream,draw=mkOrange,minimum width=17mm] (mot) {\textbf{Motor}\\ elec.$\to$mech.};
\node[mkblk,right=of mot,minimum width=17mm] (tr) {Transmission\\ gear / screw};
\node[mkblk,right=of tr,minimum width=13mm] (ld) {Load\\ $J_L,\,T_L$};
\draw[mksig,line width=1.4pt,mkOrange] (grid) -- (drv);
\draw[mksig,line width=1.4pt,mkOrange] (drv) -- node[above,mklab]{$v,\,i$} (mot);
\draw[mksig,line width=1.4pt,mkOrange] (mot) -- node[above,mklab]{$T,\,\omega$} (tr);
\draw[mksig,line width=1.4pt,mkOrange] (tr) -- node[above,mklab]{$F,\,v$} (ld);
\node[mkblk,below=8mm of mot,minimum width=15mm,fill=white] (sens) {Sensor\\ (encoder)};
\node[mkblk,below=8mm of drv,minimum width=15mm,fill=white] (ctl) {Controller\\ $C(s)$};
\draw[mksig] (mot.south) -- (sens.north);
\draw[mksig] (sens) -- (ctl);
\draw[mksig] (ctl) -- (drv);
\node[left=6mm of ctl,mklab,align=center] (ref) {reference\\ $\theta^*,\omega^*$};
\draw[mksig] (ref) -- (ctl);
% losses
\foreach \n/\t in {drv/{switching,\\ conduction},mot/{copper $i^2R$,\\ iron, friction},tr/{gear mesh,\\ friction}}{
  \draw[-{Stealth},mkRed,line width=.6pt] (\n.north) -- ++(0,5mm) node[above,font=\tiny,text=mkRed,align=center]{\t};}
\end{tikzpicture}

\vspace{2mm}
\begin{columns}[T]
\begin{column}{.5\textwidth}
\begin{eqbox}\small
$P_\text{el}=v\,i \;\xrightarrow{\text{motor}}\; P_\text{mech}=T\,\omega=P_\text{el}-P_\text{loss}$\\[1mm]
$\eta=\dfrac{T\omega}{v\,i}$, \quad typical servo motor $\eta\approx 85\text{--}95\,\%$
\end{eqbox}
\end{column}
\begin{column}{.46\textwidth}
\footnotesize
\begin{itemize}\setlength\itemsep{0pt}
\item Conversion is \hl{reversible}: motoring and generating (braking).
\item Coupling medium: the \term{magnetic field} in the air gap.
\item Losses become heat $\Rightarrow$ the \hl{thermal limit} sets continuous torque.
\end{itemize}
\end{column}
\end{columns}
\end{frame}

% ---------------------------------------------------------------------
\begin{frame}{Lorentz force and Faraday's law: the two halves of every motor}
\begin{columns}[T]
\begin{column}{.43\textwidth}
\centering
\begin{tikzpicture}[scale=.95]
% magnet poles
\fill[mkRed!70] (0,2.4) rectangle (4,2.9); \node[white,font=\scriptsize\bfseries] at (2,2.65) {N};
\fill[mkMid] (0,-.4) rectangle (4,-.9); \node[white,font=\scriptsize\bfseries] at (2,-.65) {S};
\foreach \x in {0.4,1.0,...,3.8} \draw[-{Stealth[length=1.6mm]},mkSky] (\x,2.3) -- (\x,-.3);
\node[mkSky,font=\scriptsize] at (4.25,1.0) {$\bm{B}$};
% rails & conductor
\draw[mkGray,line width=1pt] (0.2,0.3) -- (3.9,0.3); \draw[mkGray,line width=1pt] (0.2,1.7) -- (3.9,1.7);
\fill[mkOrange] (2.1,0.2) rectangle (2.35,1.8);
\draw[-{Stealth},mkOrange,line width=1pt] (2.22,0.45) -- (2.22,1.55);
\node[font=\scriptsize,text=mkOrange,right] at (2.35,1.25) {$i$};
\draw[-{Stealth},mkNavy,line width=1.4pt] (2.5,1.0) -- (3.5,1.0) node[right,font=\scriptsize]{$F,\,v$};
\draw[|-|,mkNavy] (1.8,0.3) -- node[left,font=\scriptsize]{$l$} (1.8,1.7);
\node[font=\scriptsize] at (2,-1.2) {\color{mkGray}($\bm{B}$ into the plane in the real geometry)};
\end{tikzpicture}
\end{column}
\begin{column}{.55\textwidth}
\begin{eqbox}\small
\textbf{Lorentz force} (conductor of length $l$):\quad $\bm{F}=i\,\bm{l}\times\bm{B}\;\Rightarrow\;F=B\,l\,i$\\[1mm]
\textbf{Faraday / motional EMF:}\quad $e=-\dfrac{d\lambda}{dt}\;\Rightarrow\;e=B\,l\,v$
\end{eqbox}
\vspace{1mm}
\footnotesize
\textbf{Power balance} (ideal, lossless): \quad $e\,i = (B l v)\,i = (B l i)\,v = F\,v$\\[1mm]
For a coil with $N$ turns at radius $r$ ($2N$ active sides):
\[
T = \underbrace{2NBlr}_{k_t}\, i, \qquad e=\underbrace{2NBlr}_{k_e}\,\omega
\]
\begin{keybox}[Consequence]
The same geometry gives \hl{torque per ampere} and \hl{volts per rad/s}: in SI units $k_t=k_e$
(\si{\newton\metre\per\ampere} $\equiv$ \si{\volt\second\per\radian}).
\end{keybox}
\end{column}
\end{columns}
\end{frame}

% ---------------------------------------------------------------------
\begin{frame}{Reluctance force and the co-energy method}
\begin{columns}[T]
\begin{column}{.36\textwidth}
\centering
\begin{tikzpicture}[scale=.8]
% C core
\fill[mkGray!50] (0,0) rectangle (0.6,3); \fill[mkGray!50] (0,0) rectangle (2.4,0.6); \fill[mkGray!50] (0,2.4) rectangle (2.4,3);
% coil
\foreach \y in {0.9,1.2,...,2.1} \draw[mkOrange,line width=1.6pt] (-0.15,\y) -- (0.75,\y);
\node[font=\scriptsize,text=mkOrange] at (-0.5,1.5) {$N,i$};
% rotor bar rotated
\begin{scope}[shift={(2.85,1.5)},rotate=30]
\fill[mkMid] (-0.25,-1.1) rectangle (0.25,1.1);
\end{scope}
\draw[-{Stealth},mkNavy,line width=.8pt] (3.9,0.5) arc (-40:40:1.3) node[right,font=\scriptsize]{$T$};
\node[font=\scriptsize] at (2.85,-0.3) {rotor angle $\theta$};
\end{tikzpicture}
\begin{tikzpicture}
\begin{axis}[mkplot,width=5.2cm,height=3.4cm,xmin=0,xmax=180,xtick={0,45,90,135,180},
  xlabel={$\theta$ (deg)},ymin=-1.1,ymax=2.1,legend pos=south west,legend columns=2]
\addplot[mkBlue,thick,domain=0:180,samples=60]{1+0.8*cos(2*x)}; \addlegendentry{$L(\theta)/L_0$}
\addplot[mkOrange,thick,domain=0:180,samples=60]{-0.8*sin(2*x)}; \addlegendentry{$T/T_0$}
\end{axis}
\end{tikzpicture}
\end{column}
\begin{column}{.62\textwidth}
\footnotesize
Iron is attracted towards the position of \term{minimum reluctance} (maximum inductance), even without magnets.
\begin{eqbox}\small
Co-energy: $W'(i,\theta)=\displaystyle\int_0^{i}\lambda(i',\theta)\,di'$,\qquad
$T=\left.\dfrac{\partial W'}{\partial\theta}\right|_{i=\text{const}}$\\[1mm]
Linear magnetics, $\lambda=L(\theta)\,i$:\qquad $T=\tfrac12\,i^2\,\dfrac{dL}{d\theta}$, \quad $F=\tfrac12\,i^2\,\dfrac{dL}{dx}$
\end{eqbox}
\begin{itemize}\setlength\itemsep{0pt}
\item Torque $\propto i^2$: independent of current sign $\Rightarrow$ solenoids, switched-reluctance and stepper machines.
\item $L(\theta)=L_0+L_2\cos2\theta$ gives $T=-L_2\,i^2\sin2\theta$ (plot).
\item In IPM synchronous motors a \hl{reluctance torque} adds to the magnet torque (saliency $L_q>L_d$).
\item With saturation, use the full co-energy $W'$, not $\tfrac12 Li^2$.
\end{itemize}
{\scriptsize\color{mkGray}[Fitzgerald, Kingsley \& Umans, 2003, ch.\,3]}
\end{column}
\end{columns}
\end{frame}

% ---------------------------------------------------------------------
\begin{frame}{Torque, speed and power}
\begin{columns}[T]
\begin{column}{.5\textwidth}
\begin{tikzpicture}
\begin{axis}[mkplot,width=7.2cm,height=5cm,xmin=0,xmax=2.6,ymin=0,ymax=1.25,
  xlabel={speed $\omega/\omega_\text{b}$},ylabel={normalised},legend pos=north east,
  xtick={0,0.5,1,1.5,2,2.5}]
\addplot[mkBlue,thick,domain=0:1]{1}; \addlegendentry{torque $T/T_\text{N}$}
\addplot[mkBlue,thick,domain=1:2.5,samples=40,forget plot]{1/x};
\addplot[mkOrange,thick,domain=0:1]{x}; \addlegendentry{power $P/P_\text{N}$}
\addplot[mkOrange,thick,domain=1:2.5,forget plot]{1};
\draw[dashed,mkGray] (axis cs:1,0) -- (axis cs:1,1.2);
\node[font=\tiny,text=mkMid] at (axis cs:0.5,1.12) {constant torque};
\node[font=\tiny,text=mkMid] at (axis cs:1.8,1.12) {constant power};
\node[font=\tiny,text=mkGray] at (axis cs:1.8,0.3) {(field weakening)};
\end{axis}
\end{tikzpicture}
\end{column}
\begin{column}{.48\textwidth}
\begin{eqbox}\small
$P=T\,\omega$,\qquad $\omega=\dfrac{2\pi n}{60}$ \quad ($n$ in rpm)\\[1mm]
$P\,[\si{\kilo\watt}]=\dfrac{T\,[\si{\newton\metre}]\cdot n\,[\si{rpm}]}{9549}$
\end{eqbox}
\footnotesize
\begin{itemize}\setlength\itemsep{1pt}
\item Torque is limited by \hl{current} (thermal, magnet demagnetisation).
\item Speed is limited by \hl{voltage} (back-EMF $k_e\omega$ must stay below the supply) and by mechanics.
\item Above base speed $\omega_\text{b}$, field weakening trades torque for speed at $\approx$ constant power.
\item Linear motion: $P=F\,v$; screw/belt convert $T,\omega\leftrightarrow F,v$.
\end{itemize}
\begin{notebox}
\textbf{Example:} \SI{1.27}{\newton\metre} at \SI{3000}{rpm} $\Rightarrow$ $\omega=\SI{314}{\radian\per\second}$, $P=\SI{400}{\watt}$.
\end{notebox}
\end{column}
\end{columns}
\end{frame}

% ---------------------------------------------------------------------
\begin{frame}{Classification of motors and actuators}
\centering
\resizebox{\textwidth}{!}{\begin{tikzpicture}[
  lvl1/.style={mkblk,fill=mkBlue,text=white,minimum width=24mm,font=\scriptsize\bfseries},
  lvl2/.style={mkblk,minimum width=20mm,text width=20mm},
  lvl3/.style={font=\tiny,text=mkNavy,align=left,text width=22mm},
  edge/.style={mkGray,line width=.5pt}]
\node[mkblk,fill=mkNavy,text=white,font=\small\bfseries,minimum width=40mm] (root) {Actuators};
\node[lvl1,below left=6mm and 26mm of root] (em) {Electromagnetic};
\node[lvl1,below right=6mm and 26mm of root] (oth) {Other principles};
\draw[edge] (root) -- (em); \draw[edge] (root) -- (oth);
\node[lvl2,below=7mm of em,xshift=-36mm] (dc) {DC commutator};
\node[lvl2,right=2mm of dc] (pm) {PM brushless};
\node[lvl2,right=2mm of pm] (im) {Induction};
\node[lvl2,right=2mm of im] (rel) {Reluctance / stepper};
\node[lvl2,right=2mm of rel] (lin) {Linear direct};
\foreach \n in {dc,pm,im,rel,lin} \draw[edge] (em.south) -- (\n.north);
\node[lvl3,below=1mm of dc] {PM-excited (servo)\\ series / shunt\\ coreless (ironless)};
\node[lvl3,below=1mm of pm] {BLDC (trapezoidal)\\ PMSM: SPM, IPM\\ torque motors};
\node[lvl3,below=1mm of im] {squirrel cage\\ wound rotor\\ (V/f, vector ctrl.)};
\node[lvl3,below=1mm of rel] {hybrid stepper\\ switched reluctance\\ synchronous rel.};
\node[lvl3,below=1mm of lin] {iron-core / ironless\\ linear motors\\ voice coil, solenoid};
\node[lvl2,right=8mm of lin,fill=mkCream,draw=mkOrange] (fl) {Fluid power};
\node[lvl2,right=2mm of fl,fill=mkCream,draw=mkOrange] (sm) {Smart materials};
\draw[edge] (oth.south) -- (fl.north); \draw[edge] (oth.south) -- (sm.north);
\node[lvl3,below=1mm of fl] {hydraulic cylinders\\ pneumatic cylinders};
\node[lvl3,below=1mm of sm] {piezo stacks\\ shape-memory alloy\\ magnetostrictive};
\end{tikzpicture}}

\vspace{3mm}
\footnotesize
\begin{tabular}{@{}lcccc@{}}
\toprule
 & \textbf{PM servo motor} & \textbf{Hydraulic} & \textbf{Pneumatic} & \textbf{Piezo}\\
\midrule
Force/torque density & high & \hl{very high} & medium & high (tiny stroke)\\
Bandwidth / precision & \hl{high / high} & medium / medium & low / low & \hl{very high} / nm\\
Controllability, cleanliness & \hl{excellent} & leakage, pump & compressible air & HV drive\\
\bottomrule
\end{tabular}
\end{frame}

% ---------------------------------------------------------------------
\begin{frame}{Session 1.1 -- Summary / key take-aways}
\begin{columns}[T]
\begin{column}{.55\textwidth}
\begin{itemize}
\item A motor converts $v\,i$ into $T\,\omega$ through the air-gap magnetic field; the process is reversible.
\item \term{Lorentz}: $F=Bli$ \;$\Leftrightarrow$\; \term{Faraday}: $e=Blv$ $\Rightarrow$ $k_t=k_e$ (SI).
\item \term{Reluctance} torque $T=\tfrac12 i^2\,dL/d\theta$ needs no magnets (steppers, SRM, IPM saliency).
\item Current limits torque; voltage limits speed; heat limits the continuous operating point.
\end{itemize}
\end{column}
\begin{column}{.42\textwidth}
\begin{keybox}[Remember]
\[
P=T\omega=Fv,\qquad \omega=\frac{2\pi n}{60}
\]
\[
T=k_t\,i,\qquad e=k_e\,\omega
\]
\end{keybox}
\end{column}
\end{columns}
\end{frame}

% =====================================================================
\section{1.2 Structure of electrical machines}
% =====================================================================

\begin{frame}{Anatomy of an industrial servo motor}
\begin{columns}[T]
\begin{column}{.47\textwidth}
\centering
\begin{tikzpicture}[scale=.72,lbl/.style={font=\tiny,text=mkNavy},ld/.style={mkGray,line width=.3pt}]
% housing
\draw[mkGray,line width=.8pt,fill=mkLight] (0,0) rectangle (6.4,3.4);
\foreach \x in {0.6,1.0,...,4.6} {\fill[mkGray!40] (\x,3.4) rectangle (\x+0.15,3.6); \fill[mkGray!40] (\x,-0.2) rectangle (\x+0.15,0);}
% flange (drive end)
\fill[mkGray!60] (-0.35,-0.35) rectangle (0,3.75);
% stator stack
\fill[mkSky!60] (1.0,2.3) rectangle (4.4,3.2); \fill[mkSky!60] (1.0,0.2) rectangle (4.4,1.1);
% end windings
\foreach \y in {2.3,0.2} {\fill[mkOrange!80] (0.6,\y+0.1) rectangle (1.0,\y+0.8); \fill[mkOrange!80] (4.4,\y+0.1) rectangle (4.8,\y+0.8);}
% rotor with magnets
\fill[mkGray!50] (1.0,1.25) rectangle (4.4,2.15);
\fill[mkRed!70] (1.0,2.0) rectangle (4.4,2.15); \fill[mkMid] (1.0,1.25) rectangle (4.4,1.4);
% shaft
\fill[mkGray] (-1.3,1.55) rectangle (5.9,1.85);
% bearings
\foreach \x in {0.25,5.0} {\draw[mkNavy,fill=white] (\x,1.85) rectangle (\x+0.35,2.2); \draw[mkNavy] (\x,1.85) -- (\x+0.35,2.2) (\x,2.2) -- (\x+0.35,1.85);
  \draw[mkNavy,fill=white] (\x,1.2) rectangle (\x+0.35,1.55); \draw[mkNavy] (\x,1.2) -- (\x+0.35,1.55) (\x,1.55) -- (\x+0.35,1.2);}
% brake
\fill[mkAmber!70] (5.45,1.0) rectangle (5.7,2.4);
% encoder
\fill[mkNavy] (5.9,1.1) rectangle (5.98,2.3); \fill[mkGreen!70] (6.05,2.0) rectangle (6.3,2.3);
% connector
\fill[mkGray!70] (2.6,3.6) rectangle (3.4,4.1);
% labels
\draw[ld] (2.7,2.75) -- (2.7,4.5) node[above,lbl]{stator stack + winding};
\draw[ld] (0.8,2.8) -- (0.2,4.5) node[above,lbl]{end winding};
\draw[ld] (3.4,3.85) -- (4.5,4.5) node[above,lbl]{connector};
\draw[ld] (2.0,2.07) -- (1.6,-0.9) node[below,lbl]{PM rotor};
\draw[ld] (0.42,1.3) -- (0.2,-0.9) node[below,lbl]{bearing};
\draw[ld] (-0.9,1.7) -- (-1.1,-0.9) node[below,lbl]{shaft (DE)};
\draw[ld] (5.57,1.1) -- (4.4,-0.9) node[below,lbl]{brake};
\draw[ld] (6.1,2.1) -- (6.6,-0.9) node[below,lbl]{encoder (NDE)};
\draw[ld] (3.0,0.9) -- (3.0,-0.9) node[below,lbl]{air gap};
\end{tikzpicture}\\
{\tiny\color{mkGray} Schematic longitudinal section of a PM servo motor (not to scale).}
\end{column}
\begin{column}{.53\textwidth}
\footnotesize
\begin{itemize}\setlength\itemsep{2pt}
\item \term{Stator}: laminated iron stack with slots and (three-phase) copper windings -- the heat source is on the outside, close to the housing.
\item \term{Rotor}: shaft + magnets (PM) or cage (induction); low inertia is a design goal for servo motors (long, slim rotors).
\item \term{Air gap}: $\approx$ \SIrange{0.3}{1}{\milli\metre}; all torque is produced here.
\item \term{Feedback device}: encoder or resolver on the non-drive end (Day 5).
\item \term{Holding brake} (optional, spring-applied), \term{bearings}, \term{shaft seal}, \term{connectors}.
\item Ratings defined per IEC 60034-1 (duty types S1--S10), protection class IP65/67 typical.
\end{itemize}
\begin{notebox}
Industrial ``AC servo'' usually means a \hl{PMSM} with sinusoidal back-EMF, encoder and a matched drive.
\end{notebox}
\end{column}
\end{columns}
\end{frame}

% ---------------------------------------------------------------------
\begin{frame}{Brushed (commutator) DC motor: construction}
\begin{columns}[T]
\begin{column}{.40\textwidth}
\centering
\begin{tikzpicture}[scale=.9]
\fill[mkGray!40] (0,0) circle (2.0); \fill[white] (0,0) circle (1.65);
% magnets
\fill[mkRed!70] (100:1.65) arc (100:260:1.65) -- (260:1.35) arc (260:100:1.35) -- cycle;
\fill[mkMid] (-80:1.65) arc (-80:80:1.65) -- (80:1.35) arc (80:-80:1.35) -- cycle;
\node[white,font=\scriptsize\bfseries] at (180:1.5) {N}; \node[white,font=\scriptsize\bfseries] at (0:1.5) {S};
% rotor
\fill[mkSky!50] (0,0) circle (1.22);
\foreach \a in {0,30,...,330}{ \fill[white] (\a:1.05) circle (0.1);}
\foreach \a in {120,150,180,210,240}{ \draw[mkOrange,line width=.7pt] (\a:1.05) circle (0.09); \fill[mkOrange] (\a:1.05) circle (0.025);}
\foreach \a in {-60,-30,0,30,60}{ \draw[mkOrange,line width=.7pt] (\a:1.05) circle (0.09);
   \draw[mkOrange,line width=.6pt] ($(\a:1.05)+(-0.06,-0.06)$) -- ($(\a:1.05)+(0.06,0.06)$) ($(\a:1.05)+(-0.06,0.06)$) -- ($(\a:1.05)+(0.06,-0.06)$);}
% commutator & brushes
\fill[mkAmber] (0,0) circle (0.42); \foreach \a in {0,45,...,315} \draw[white] (0,0) -- (\a:0.42);
\fill[mkNavy] (-0.42,-0.1) rectangle (-0.7,0.1); \fill[mkNavy] (0.42,-0.1) rectangle (0.7,0.1);
\fill[mkGray] (0,0) circle (0.12);
\draw[-{Stealth},mkNavy,line width=1pt] (60:2.25) arc (60:120:2.25) node[left,font=\scriptsize]{$T$};
\node[font=\tiny,text=mkGray] at (0,-2.3) {stator magnets, armature, commutator, brushes};
\end{tikzpicture}
\end{column}
\begin{column}{.30\textwidth}
\centering
\figph{width=\linewidth,height=4.6cm,keepaspectratio}{d1_brushed_dc.jpg}\\
\credit{Jfer91, CC BY-SA 3.0, via Wikimedia Commons}
\end{column}
\begin{column}{.28\textwidth}
\footnotesize
\begin{itemize}\setlength\itemsep{2pt}
\item Field from \term{stator magnets} (or field winding).
\item Current-carrying \term{armature} rotates.
\item \term{Commutator + brushes} reverse coil currents so that conductors under N always carry the same direction.
\item Armature MMF stays \hl{perpendicular} to the field $\Rightarrow$ maximum torque per ampere, automatically.
\end{itemize}
\end{column}
\end{columns}
\end{frame}

% ---------------------------------------------------------------------
\begin{frame}{Mechanical commutation and torque ripple}
\begin{columns}[T]
\begin{column}{.52\textwidth}
\begin{tikzpicture}
\begin{axis}[mkplot,width=7.4cm,height=4.2cm,xmin=0,xmax=360,ymin=0,ymax=1.15,
  xtick={0,90,180,270,360},xlabel={rotor angle $\theta$ (deg)},ylabel={$T/T_\text{max}$},legend pos=south east,legend columns=1]
\addplot[mkSky,thick,domain=0:360,samples=181]{abs(sin(x))}; \addlegendentry{1 coil, 2 segments}
\addplot[mkOrange,thick,domain=0:360,samples=361]{max(abs(sin(x)),abs(sin(x+60)),abs(sin(x+120)))}; \addlegendentry{3 coils}
\addplot[mkBlue,thick,domain=0:360,samples=721]{max(abs(sin(x)),abs(sin(x+15)),abs(sin(x+30)),abs(sin(x+45)),abs(sin(x+60)),abs(sin(x+75)),abs(sin(x+90)),abs(sin(x+105)),abs(sin(x+120)),abs(sin(x+135)),abs(sin(x+150)),abs(sin(x+165)))}; \addlegendentry{12 coils}
\end{axis}
\end{tikzpicture}\\
{\tiny\color{mkGray} Idealised: only the best-aligned coil group is counted.}
\end{column}
\begin{column}{.46\textwidth}
\footnotesize
\begin{itemize}\setlength\itemsep{2pt}
\item The commutator is a \hl{mechanical inverter}: it switches each coil's current as it passes the neutral zone.
\item With $n$ coil groups, peak-to-peak ripple $\approx 1-\cos(\pi/2n)$:\\
$n=3$: 13\,\%, \quad $n=12$: 0.9\,\%.
\item Commutation sparks, brush wear (typically a few thousand hours), EMI, speed limit from brush bounce.
\item Heat is generated in the \hl{rotor} $\Rightarrow$ poor cooling path, low continuous torque per volume.
\item Simple drive: one H-bridge.
\end{itemize}
\begin{keybox}[Idea for BLDC/PMSM]
Replace the commutator by \hl{power electronics + rotor-position sensor} $\Rightarrow$ electronic commutation.
\end{keybox}
\end{column}
\end{columns}
\end{frame}

% ---------------------------------------------------------------------
\begin{frame}{Permanent-magnet brushless machines: BLDC and PMSM}
\begin{columns}[T]
\begin{column}{.52\textwidth}
\centering
\begin{tikzpicture}[scale=.75]
\foreach \cx/\name in {0/SPM, 4.6/IPM}{
\begin{scope}[shift={(\cx,0)}]
 \fill[mkGray!40] (0,0) circle (2.0); \fill[white] (0,0) circle (1.48);
 \foreach \a in {0,30,...,330}{ \fill[mkGray!40,rotate=\a] (1.4,-0.13) rectangle (1.55,0.13);}
 \foreach \a in {15,45,...,345}{ \fill[mkOrange!60,rotate=\a] (1.55,-0.1) rectangle (1.85,0.1);}
 \fill[mkSky!50] (0,0) circle (1.1);
 \fill[mkGray] (0,0) circle (0.25);
 \node[font=\scriptsize\bfseries,text=mkNavy] at (0,-2.4) {\name};
\end{scope}}
% SPM magnets on surface
\foreach \a/\c in {0/mkRed!70,90/mkMid,180/mkRed!70,270/mkMid}{
  \fill[\c] (\a-38:1.32) arc (\a-38:\a+38:1.32) -- (\a+38:1.1) arc (\a+38:\a-38:1.1) -- cycle;}
% IPM V-shaped buried magnets
\begin{scope}[shift={(4.6,0)}]
\foreach \a/\c in {0/mkRed!70,90/mkMid,180/mkRed!70,270/mkMid}{
  \begin{scope}[rotate=\a]
  \fill[\c,rotate around={25:(0.75,0.25)}] (0.45,0.2) rectangle (1.0,0.32);
  \fill[\c,rotate around={-25:(0.75,-0.25)}] (0.45,-0.32) rectangle (1.0,-0.2);
  \end{scope}}
\end{scope}
\end{tikzpicture}\\[1mm]
\begin{tabular}{@{}ll@{}}
\footnotesize\textbf{SPM}: & \scriptsize magnets glued on surface, $L_d\approx L_q$\\
\footnotesize\textbf{IPM}: & \scriptsize buried magnets, $L_q>L_d$, reluctance torque
\end{tabular}
\end{column}
\begin{column}{.46\textwidth}
\footnotesize
\begin{itemize}\setlength\itemsep{2pt}
\item ``Inside-out'' DC motor: magnets on the \term{rotor}, three-phase winding in the \term{stator} $\Rightarrow$ heat close to the housing, no brushes.
\item \term{BLDC}: concentrated windings, \hl{trapezoidal} back-EMF, six-step (block) commutation.
\item \term{PMSM}: distributed (or optimised) windings, \hl{sinusoidal} back-EMF, sinusoidal / field-oriented control.
\item $p$ pole pairs: $\theta_e=p\,\theta_m$, \; $\omega_e=p\,\omega_m$.
\item Servo: mostly SPM; traction/spindles: IPM.
\end{itemize}
\centering
\figph{width=.6\linewidth,height=1.6cm,keepaspectratio}{d1_bldc.jpg}\\
\credit{Sebastian Koppehel, CC BY 3.0, via Wikimedia Commons}
\end{column}
\end{columns}
\end{frame}

% ---------------------------------------------------------------------
\begin{frame}{Electronic commutation I: six-step (block) with Hall sensors}
\begin{columns}[T]
\begin{column}{.56\textwidth}
\begin{tikzpicture}[x=0.0155cm,y=0.36cm]
\foreach \y/\n in {9/$H_A$,7.5/$H_B$,6/$H_C$} \node[font=\tiny,anchor=east] at (-5,\y+0.4) {\n};
\draw[mkBlue,thick] (0,9) -- (30,9) -- (30,9.8) -- (210,9.8) -- (210,9) -- (360,9);
\draw[mkBlue,thick] (0,7.5) -- (150,7.5) -- (150,8.3) -- (330,8.3) -- (330,7.5) -- (360,7.5);
\draw[mkBlue,thick] (0,6.8) -- (90,6.8) -- (90,6) -- (270,6) -- (270,6.8) -- (360,6.8);
% back-emf & current a
\node[font=\tiny,anchor=east] at (-5,3) {$e_a$};
\draw[mkGray!60] (0,3) -- (360,3);
\draw[mkSky,thick] (0,3) -- (30,4.5) -- (150,4.5) -- (210,1.5) -- (330,1.5) -- (360,3);
\node[font=\tiny,anchor=east] at (-5,0) {$i_a$};
\draw[mkGray!60] (0,0) -- (360,0);
\draw[mkOrange,thick] (0,0) -- (30,0) -- (30,1.2) -- (150,1.2) -- (150,0) -- (210,0) -- (210,-1.2) -- (330,-1.2) -- (330,0) -- (360,0);
\foreach \x in {30,90,...,330} \draw[mkGray,dashed,very thin] (\x,-1.6) -- (\x,10);
\foreach \x/\l in {60/A+B$-$,120/A+C$-$,180/B+C$-$,240/B+A$-$,300/C+A$-$} \node[font=\tiny,rotate=90,text=mkMid] at (\x,-2.8) {\l};
\node[font=\tiny,rotate=90,text=mkMid] at (15,-2.8) {C+B$-$};
\node[font=\tiny,rotate=90,text=mkMid] at (345,-2.8) {C+B$-$};
\draw[-{Stealth}] (0,-4.6) -- (370,-4.6) node[right,font=\tiny]{$\theta_e$};
\foreach \x in {0,60,...,360} \node[font=\tiny,below] at (\x,-4.6) {\x$^\circ$};
\end{tikzpicture}
\end{column}
\begin{column}{.42\textwidth}
\scriptsize
\renewcommand{\arraystretch}{1.0}
\begin{tabular}{@{}cccl@{}}
\toprule
$\theta_e$ & $H_AH_BH_C$ & on & current path\\
\midrule
30--90 & 101 & A+ B$-$ & $i_a=-i_b=I$\\
90--150 & 100 & A+ C$-$ & $i_a=-i_c=I$\\
150--210 & 110 & B+ C$-$ & $i_b=-i_c=I$\\
210--270 & 010 & B+ A$-$ & $i_b=-i_a=I$\\
270--330 & 011 & C+ A$-$ & $i_c=-i_a=I$\\
330--30 & 001 & C+ B$-$ & $i_c=-i_b=I$\\
\bottomrule
\end{tabular}

\vspace{1mm}
\footnotesize
\begin{itemize}\setlength\itemsep{1pt}
\item Three Hall sensors give \hl{60$^\circ_e$ resolution}: 6 states per electrical period.
\item Two phases conduct, one floats: $T=\dfrac{e_ai_a+e_bi_b+e_ci_c}{\omega_m}=\dfrac{2E\,I}{\omega_m}$.
\item Current transfers at each step $\Rightarrow$ \hl{commutation torque ripple} (10--20\,\%).
\item Sensorless option: detect back-EMF zero crossing in the floating phase.
\end{itemize}
{\tiny\color{mkGray} Hall mounting/coding differs between manufacturers -- table is for the alignment shown.}
\end{column}
\end{columns}
\end{frame}

% ---------------------------------------------------------------------
\begin{frame}{Electronic commutation II: sinusoidal commutation (PMSM)}
\begin{columns}[T]
\begin{column}{.4\textwidth}
\centering
\begin{tikzpicture}[scale=1.15]
\draw[mkGray!60] (0,0) circle (1.5);
\foreach \a/\n in {0/a,120/b,240/c}{\draw[mkGray,-{Stealth}] (0,0) -- (\a:1.75) node[font=\scriptsize,text=mkGray] at (\a:1.95) {\n};}
\draw[-{Stealth[length=2.5mm]},mkMid,line width=1.6pt] (0,0) -- (35:1.2) node[pos=1.1,font=\scriptsize,text=mkMid,right]{$d$: $\psi_\text{PM}$};
\draw[-{Stealth[length=2.5mm]},mkOrange,line width=1.6pt] (0,0) -- (125:1.3) node[above left,font=\scriptsize,text=mkOrange]{$\bm{i}_s=i_q$};
\draw[mkGray,dashed] (0,0) -- (125:1.6) node[above,font=\tiny] {$q$};
\coordinate (O) at (0,0); \coordinate (X) at (1,0); \coordinate (D) at (35:1); \coordinate (Q) at (125:1);
\pic[draw=mkNavy,"$\theta_e$",angle radius=5mm,angle eccentricity=1.6,font=\tiny] {angle=X--O--D};
\pic[draw=mkNavy,angle radius=3.5mm,"$90^\circ$",angle eccentricity=1.9,font=\tiny] {angle=D--O--Q};
\draw[-{Stealth},mkNavy] (-20:1.65) arc (-20:15:1.65) node[right,font=\tiny]{$\omega_e$};
\end{tikzpicture}
\begin{tikzpicture}
\begin{axis}[mkplot,width=5.6cm,height=3.0cm,xmin=0,xmax=360,ymin=-1.2,ymax=1.2,xtick={0,120,240,360},
xlabel={$\theta_e$ (deg)},ytick={-1,0,1}]
\addplot[mkBlue,thick,domain=0:360,samples=90]{sin(x)};
\addplot[mkOrange,thick,domain=0:360,samples=90]{sin(x-120)};
\addplot[mkSky,thick,domain=0:360,samples=90]{sin(x+120)};
\end{axis}
\end{tikzpicture}
\end{column}
\begin{column}{.58\textwidth}
\footnotesize
Impose phase currents \hl{in phase with the sinusoidal back-EMF}:
\begin{eqbox}\small
$e_x = k_{e,\text{ph}}\,\omega_m\sin(\theta_e-\varphi_x)$, \; $i_x=\hat I\sin(\theta_e-\varphi_x)$\\ $\varphi_x\in\{0,\tfrac{2\pi}{3},\tfrac{4\pi}{3}\}$ for $x=a,b,c$\\[1mm]
$T=\dfrac{\sum_x e_x i_x}{\omega_m}=k_{e,\text{ph}}\hat I\sum_x\sin^2(\theta_e-\varphi_x)=\tfrac32\,k_{e,\text{ph}}\,\hat I=\text{const.}$
\end{eqbox}
\begin{itemize}\setlength\itemsep{1pt}
\item Needs \hl{continuous rotor angle} (encoder/resolver), not just 6 Hall states.
\item Current vector kept at $90^\circ_e$ to the magnet flux: $i_d=0$, $T=\tfrac32\,p\,\psi_\text{PM}\,i_q$.
\item Ideally \hl{ripple-free} torque ($<2\,\%$ vs.\ 10--20\,\% for six-step).
\item Implemented via Clarke/Park transforms = \term{field-oriented control} (Day 3--4).
\end{itemize}

\end{column}
\end{columns}
\end{frame}

% ---------------------------------------------------------------------
\begin{frame}{Induction (asynchronous) motor: short introduction}
\begin{columns}[T]
\begin{column}{.30\textwidth}
\centering
\figph{width=\linewidth,height=3.4cm,keepaspectratio}{d1_squirrel_cage.jpg}\\
\credit{Zureks, CC BY-SA 3.0, via Wikimedia Commons}
\vspace{1mm}
\footnotesize
\begin{itemize}\setlength\itemsep{1pt}
\item Rotating stator field induces rotor currents in a \term{squirrel cage}.
\item Rotor must \hl{slip} behind the field to see induced voltage.
\end{itemize}
\end{column}
\begin{column}{.40\textwidth}
\begin{tikzpicture}
\begin{axis}[mkplot,width=6.2cm,height=5.0cm,xmin=0,xmax=1.05,ymin=0,ymax=2.8,
  xlabel={speed $n/n_s$},ylabel={$T/T_\text{N}$}]
\addplot[mkBlue,thick,domain=0:0.999,samples=200]{2*2.5/(((1-x)/0.2)+(0.2/(1-x)))};
\addplot[mkOrange,thick,dashed,domain=0:1]{1};
\node[font=\tiny,text=mkBlue] at (axis cs:0.75,2.6) {$T_k$ (breakdown)};
\node[font=\tiny,text=mkOrange] at (axis cs:0.25,1.15) {rated load};
\node[circle,fill=mkOrange,inner sep=1.3pt] at (axis cs:0.958,1) {};
\node[font=\tiny,anchor=east] at (axis cs:0.95,0.82) {$s_\text{N}\approx4\,\%$};
\end{axis}
\end{tikzpicture}
\end{column}
\begin{column}{.28\textwidth}
\begin{eqbox}\footnotesize
$n_s=\dfrac{60\,f}{p}$, \quad $s=\dfrac{n_s-n}{n_s}$\\[1mm]
Kloss: $\dfrac{T}{T_k}\approx\dfrac{2}{s/s_k+s_k/s}$
\end{eqbox}
\footnotesize
\begin{itemize}\setlength\itemsep{1pt}
\item Robust, cheap, no magnets, field weakening easy.
\item Rotor losses $\propto s$ heat the rotor.
\item Servo use: spindles, large axes, with \hl{vector control} + encoder.
\item Plot: $s_k=0.2$, $T_k=2.5\,T_N$.
\end{itemize}
\end{column}
\end{columns}
\end{frame}

% ---------------------------------------------------------------------
\begin{frame}{Stepper motor: short introduction}
\begin{columns}[T]
\begin{column}{.38\textwidth}
\centering
\figph{width=\linewidth,height=4.2cm,keepaspectratio}{d1_stepper.jpg}\\
\credit{tony\_duell, CC BY 2.0, via Wikimedia Commons}
\end{column}
\begin{column}{.30\textwidth}
\begin{tikzpicture}
\begin{axis}[mkplot,width=4.4cm,height=4.4cm,xmin=0,xmax=3000,ymin=0,ymax=1.1,
  xlabel={step rate (full steps/s)},ylabel={$T/T_\text{hold}$},xtick={0,1000,2000,3000},
  scaled x ticks=false]
\addplot[mkBlue,thick,domain=0:3000,samples=80]{0.95/(1+(x/900)^1.6)};
\addplot[mkSky,thick,dashed,domain=0:1500,samples=60]{0.95/(1+(x/450)^1.6)};
\node[font=\tiny,text=mkBlue] at (axis cs:1900,0.38) {pull-out};
\node[font=\tiny,text=mkSky] at (axis cs:600,0.12) {pull-in};
\end{axis}
\end{tikzpicture}\\
{\tiny\color{mkGray}qualitative shape}
\end{column}
\begin{column}{.30\textwidth}
\footnotesize
\begin{itemize}\setlength\itemsep{1pt}
\item Hybrid stepper: toothed rotor with axial magnet, 2 phases.
\item Step angle: $\theta_\text{step}=\dfrac{360^\circ}{2\,m\,N_r}$;\\ $m=2$, $N_r=50$ $\Rightarrow$ \hl{1.8$^\circ$} (200 steps/rev).
\item \term{Open-loop} positioning by counting steps; micro-stepping for smoothness.
\item Torque falls with speed ($L$, back-EMF).
\item Risk: \hl{lost steps} if load torque exceeds pull-out torque; resonance at low/medium speed.
\end{itemize}
\end{column}
\end{columns}
\end{frame}

% ---------------------------------------------------------------------
\begin{frame}{Comparison of machine types for motion control}
\centering
\footnotesize
\renewcommand{\arraystretch}{1.1}
\begin{tabularx}{\textwidth}{@{}l*{5}{>{\centering\arraybackslash}X}@{}}
\toprule
 & \textbf{Brushed DC} & \textbf{BLDC} & \textbf{PMSM (AC servo)} & \textbf{Induction} & \textbf{Hybrid stepper}\\
\midrule
Commutation & mechanical & 6-step, Hall & sinusoidal / FOC & slip, vector ctrl. & phase sequencing\\
Torque density & medium & high & \textcolor{mkGreen}{\textbf{very high}} & medium & high at low speed\\
Torque ripple & low--medium & \textcolor{mkRed}{medium--high} & \textcolor{mkGreen}{\textbf{very low}} & low & \textcolor{mkRed}{high (detent)}\\
Speed range & limited (brushes) & high & high & \textcolor{mkGreen}{\textbf{very high}} & \textcolor{mkRed}{low}\\
Feedback & optional & Hall & encoder/resolver & encoder (servo) & \textcolor{mkGreen}{none needed}\\
Maintenance & \textcolor{mkRed}{brushes} & none & none & none & none\\
Drive complexity & \textcolor{mkGreen}{\textbf{H-bridge}} & 3-ph. inverter & 3-ph. inverter + FOC & 3-ph. inverter & 2$\times$H-bridge\\
Use & toys, small tools & fans, pumps, drones & robots, machine tools & spindles, pumps & printers, 3D printers\\
\bottomrule
\end{tabularx}

\begin{notebox}
Rule of thumb for industrial servo axes: \hl{PMSM} below $\approx$\SI{10}{\kilo\watt} and where dynamics/precision dominate; induction for large spindles and rugged, high-speed applications.
\end{notebox}
\end{frame}

% ---------------------------------------------------------------------
\begin{frame}{Session 1.2 -- Summary / key take-aways}
\begin{columns}[T]
\begin{column}{.56\textwidth}
\begin{itemize}
\item All machines: stator, rotor, air gap; torque comes from interaction of two magnetic fields (or reluctance).
\item \term{Commutation} keeps stator and rotor fields $\approx 90^\circ$ apart: commutator (DC), six-step + Hall (BLDC), continuous angle + sinusoidal currents (PMSM).
\item BLDC/PMSM: windings on the stator $\Rightarrow$ better cooling, no wear.
\item Induction: needs slip; stepper: open-loop steps, risk of step loss.
\end{itemize}
\end{column}
\begin{column}{.40\textwidth}
\begin{keybox}[Remember]
\[\theta_e=p\,\theta_m\]
\[T_\text{PMSM}=\tfrac32\,p\,\psi_\text{PM}\,i_q\]
\[n_s=60f/p,\quad s=(n_s-n)/n_s\]
\[\theta_\text{step}=360^\circ/(2mN_r)\]
\end{keybox}
\end{column}
\end{columns}
\end{frame}

% =====================================================================
\section{1.3 Electrical and mechanical models}
% =====================================================================

\begin{frame}{Equivalent circuit of the PM DC motor}
\begin{columns}[T]
\begin{column}{.55\textwidth}
\centering
\begin{circuitikz}[scale=.64,transform shape,font=\small]
\draw (0,0) to[V,v<=$v$,color=mkNavy] (0,3)
  to[R,l=$R$,i>^=$i$,color=mkNavy] (2.6,3)
  to[L,l=$L$,color=mkNavy] (4.8,3)
  to[short] (5.6,3)
  to[sV,v^<={$e=k_e\omega$},color=mkOrange] (5.6,0) -- (0,0);
\draw[mkGray,dashed] (7.0,-0.4) -- (7.0,3.4);
% mechanical
\begin{scope}[shift={(8.3,1.5)}]
\fill[mkSky!60] (0,-0.9) rectangle (0.5,0.9);
\node[font=\scriptsize] at (0.25,-1.15) {$J$};
\draw[mkGray,line width=2pt] (0.5,0) -- (1.9,0);
\draw[mkGray,line width=2pt] (-0.7,0) -- (0,0);
\draw[-{Stealth},mkOrange,line width=.9pt] (-0.4,0.6) arc (120:240:0.65);
\node[font=\scriptsize,text=mkOrange] at (-0.7,0.9) {$T=k_t i$};
\draw[-{Stealth},mkNavy,line width=.9pt] (1.5,-0.6) arc (-120:120:0.65);
\node[font=\scriptsize] at (2.0,0.95) {$T_L$};
\node[font=\scriptsize] at (1.2,-1.0) {$\omega,\;B\omega$};
\end{scope}
\end{circuitikz}
\end{column}
\begin{column}{.43\textwidth}
\footnotesize
\begin{tabular}{@{}lll@{}}
\toprule
\textbf{Sym.} & \textbf{Meaning} & \textbf{Unit}\\
\midrule
$R$ & winding (terminal) resistance & \si{\ohm}\\
$L$ & winding inductance & \si{\henry}\\
$k_e$ & back-EMF constant & \si{\volt\second\per\radian}\\
$k_t$ & torque constant & \si{\newton\metre\per\ampere}\\
$J$ & rotor (+ load) inertia & \si{\kilo\gram\metre\squared}\\
$B$ & viscous friction coeff. & \si{\newton\metre\second\per\radian}\\
$T_L$ & load torque & \si{\newton\metre}\\
\bottomrule
\end{tabular}
\end{column}
\end{columns}
\vspace{2mm}
\begin{notebox}
\textbf{Sign convention:} motor (consumer) convention -- $v$ and $i$ positive into the terminals; $\omega$ and $T$ positive in the same direction; $T_L$ opposes motion. Valid for brushed DC and, per $q$-axis, for PMSM under FOC (Day 4).
\end{notebox}
\end{frame}

% ---------------------------------------------------------------------
\begin{frame}{Fundamental motor equations and why $k_t=k_e$}
\begin{columns}[T]
\begin{column}{.52\textwidth}
\begin{eqbox}\small
\begin{tabular}{@{}ll@{}}
\textbf{KVL:} & $v = R\,i + L\dfrac{di}{dt} + k_e\,\omega$\\[2mm]
\textbf{Newton:} & $J\dfrac{d\omega}{dt} = k_t\,i - B\,\omega - T_L$\\[2mm]
\textbf{Kinematics:} & $\dfrac{d\theta}{dt}=\omega$
\end{tabular}
\end{eqbox}
\footnotesize
\vspace{1mm}
Multiply KVL by $i$ and Newton by $\omega$:
\[
\underbrace{v\,i}_{P_\text{el}} = \underbrace{R i^2}_{\text{Cu loss}} + \underbrace{\tfrac{d}{dt}\big(\tfrac12 L i^2\big)}_{\text{magn. energy}} + \underbrace{k_e\,\omega\, i}_{P_\text{conv}}
\]
\[
\underbrace{k_t\, i\,\omega}_{P_\text{conv}} = \underbrace{\tfrac{d}{dt}\big(\tfrac12 J\omega^2\big)}_{\text{kin. energy}} + \underbrace{B\omega^2}_{\text{friction}} + \underbrace{T_L\,\omega}_{\text{to load}}
\]
\end{column}
\begin{column}{.45\textwidth}
\begin{keybox}[Energy conservation $\Rightarrow k_t{=}k_e$]
The converted power must be identical on both sides: $e\,i = T\,\omega$\\
$\Rightarrow\; k_e\,\omega\,i = k_t\,i\,\omega \;\Rightarrow\; \bm{k_t = k_e}$ \;(SI units).
\end{keybox}
\footnotesize
\begin{itemize}\setlength\itemsep{1pt}
\item Datasheets often give $k_e$ in \si{\volt}/krpm: \\$k_e[\si{\volt\second\per\radian}]=k_e[\si{\volt}/\text{krpm}]\,/\,104.7$.
\item Example: \SI{10.47}{\volt}/krpm $\Rightarrow$ \SI{0.1}{\volt\second\per\radian} $\Rightarrow k_t=\SI{0.1}{\newton\metre\per\ampere}$.
\item Two linear, coupled 1st-order ODEs $\Rightarrow$ a \hl{2nd-order} system from $v$ to $\omega$.
\item $k_t$ falls slightly at high current (saturation) and with magnet temperature ($\approx -0.1\,\%/\mathrm{K}$ for NdFeB).
\end{itemize}
\end{column}
\end{columns}
\end{frame}

% ---------------------------------------------------------------------
\begin{frame}{Friction and load torque}
\begin{columns}[T]
\begin{column}{.48\textwidth}
\begin{tikzpicture}
\begin{axis}[mkplot,width=7cm,height=5.2cm,xmin=-3,xmax=3,ymin=-1.9,ymax=1.9,
  axis lines=middle,xlabel={$\omega$},ylabel={$T_f$},xtick=\empty,ytick=\empty,
  every axis x label/.style={at={(current axis.right of origin)},anchor=west,font=\scriptsize},
  every axis y label/.style={at={(current axis.above origin)},anchor=south,font=\scriptsize},
  legend style={at={(0.02,0.98)},anchor=north west}]
\addplot[mkSky,thick,dashed,domain=0.001:3,samples=40]{0.25*x}; \addlegendentry{viscous $B\omega$}
\addplot[mkMid,thick,dotted,domain=0.001:3,samples=2]{0.8}; \addlegendentry{Coulomb $T_C\,\mathrm{sgn}\,\omega$}
\addplot[mkOrange,thick,domain=0.001:3,samples=120]{0.8+0.6*exp(-(x/0.35)^2)+0.25*x}; \addlegendentry{Stribeck + Coulomb + viscous}
\addplot[mkSky,thick,dashed,domain=-3:-0.001,samples=40,forget plot]{0.25*x};
\addplot[mkMid,thick,dotted,domain=-3:-0.001,samples=2,forget plot]{-0.8};
\addplot[mkOrange,thick,domain=-3:-0.001,samples=120,forget plot]{-0.8-0.6*exp(-(x/0.35)^2)+0.25*x};
\draw[mkOrange,thick] (axis cs:0,-1.4) -- (axis cs:0,1.4);
\node[font=\tiny,anchor=west,text=mkOrange] at (axis cs:0.05,1.48) {static $T_s$ (stiction)};
\end{axis}
\end{tikzpicture}
\end{column}
\begin{column}{.5\textwidth}
\begin{eqbox}\small
$T_f(\omega) = \big[T_C + (T_s-T_C)\,e^{-(\omega/\omega_s)^2}\big]\,\mathrm{sgn}(\omega) + B\,\omega$
\end{eqbox}
\footnotesize
\begin{itemize}\setlength\itemsep{1pt}
\item \term{Viscous} $B\omega$: linear, enters the transfer function.
\item \term{Coulomb} $T_C$ and \term{stiction} $T_s$: \hl{nonlinear} -- cause stick--slip, limit cycles and position errors at reversal (Day 2 linearisation, Day 4 compensation).
\item \term{Load torque} $T_L$: process force, gravity ($m g r$ on vertical axes -- constant!), spring forces, cutting forces.
\item In the linear model: $T_L$ is a \hl{disturbance input}, $B$ lumps all speed-proportional losses.
\end{itemize}
{\scriptsize\color{mkGray}[Armstrong-H\'elouvry, Dupont \& Canudas de Wit, 1994]}
\end{column}
\end{columns}
\end{frame}

% ---------------------------------------------------------------------
\begin{frame}{Steady-state characteristics at constant voltage}
\begin{columns}[T]
\begin{column}{.62\textwidth}
\begin{tikzpicture}
\begin{groupplot}[group style={group size=2 by 1,horizontal sep=12mm},mkplot,width=4.7cm,height=4.6cm,
  xmin=0,xmax=4.8,xlabel={$T$ (\si{\newton\metre})}]
\nextgroupplot[ymin=0,ymax=5000,ylabel={speed (rpm), current$\times$100 (A)},legend pos=north east,
  scaled y ticks=false,yticklabel style={/pgf/number format/fixed}]
\addplot[mkBlue,thick,domain=0:4.75]{(475-100*x)*9.549}; \addlegendentry{$n$}
\addplot[mkOrange,thick,domain=0:4.75]{(10*x+0.5)*100}; \addlegendentry{$i$}
\nextgroupplot[ymin=0,ymax=1,ylabel={$P/P_\text{max}$, $\eta$},legend pos=north east]
\addplot[mkBlue,thick,domain=0:4.75,samples=60]{x*(475-100*x)/564}; \addlegendentry{$P$}
\addplot[mkOrange,thick,domain=0.001:4.75,samples=120]{x*(475-100*x)/(48*(10*x+0.5))}; \addlegendentry{$\eta$}
\end{groupplot}
\end{tikzpicture}\\
{\tiny\color{mkGray} Example motor: $v=\SI{48}{\volt}$, $R=\SI{1}{\ohm}$, $k_t=k_e=0.1$ SI, no-load friction $T_0=\SI{0.05}{\newton\metre}$.}
\end{column}
\begin{column}{.36\textwidth}
\footnotesize
Set $d/dt=0$:
\begin{eqbox}\small
$\omega=\dfrac{v}{k_e}-\dfrac{R}{k_ek_t}\,T$
\end{eqbox}
\begin{itemize}\setlength\itemsep{1pt}
\item No-load speed $\omega_0=v/k_e$ (\SI{4584}{rpm} ideal).
\item Stall torque $T_\text{st}=k_t v/R$ (\SI{4.8}{\newton\metre}, \SI{48}{\ampere}!).
\item Slope $-R/(k_ek_t)$: speed \hl{droop} under load.
\item $P_\text{max}=v^2/(4R)$ at $T_\text{st}/2$ -- never a servo operating point (50\,\% efficiency).
\item $\eta_\text{max}=\big(1-\sqrt{i_0/i_\text{st}}\big)^2\approx 81\,\%$ near light load.
\end{itemize}
\end{column}
\end{columns}
\end{frame}

% ---------------------------------------------------------------------
\begin{frame}{Block diagram of the DC motor}
\centering
\begin{tikzpicture}[node distance=7mm]
\node (v) {$V(s)$};
\node[mksum,right=6mm of v] (s1) {};
\node[mkblk,right=of s1] (el) {$\dfrac{1}{Ls+R}$};
\node[mkblk,right=of el] (kt) {$k_t$};
\node[mksum,right=of kt] (s2) {};
\node[mkblk,right=of s2] (me) {$\dfrac{1}{Js+B}$};
\node[mkblk,right=10mm of me] (int) {$\dfrac{1}{s}$};
\node[right=6mm of int] (th) {$\Theta(s)$};
\node[above=6mm of s2] (tl) {$T_L(s)$};
\node[mkblk,below=9mm of kt,fill=mkCream,draw=mkOrange] (ke) {$k_e$};
\draw[mksig] (v) -- node[pos=.85,above,font=\tiny]{$+$} (s1);
\draw[mksig] (s1) -- (el);
\draw[mksig] (el) -- node[above,mklab]{$I$} (kt);
\draw[mksig] (kt) -- node[above,mklab]{$T_m$} node[pos=.85,above,font=\tiny]{$+$} (s2);
\draw[mksig] (tl) -- node[pos=.8,right,font=\tiny]{$-$} (s2);
\draw[mksig] (s2) -- (me);
\draw[mksig] (me) -- node[above,mklab]{$\Omega$} coordinate[pos=.4] (wtap) (int);
\draw[mksig] (int) -- (th);
\draw[mksig] (wtap) |- (ke);
\draw[mksig] (ke) -| node[pos=.9,left,font=\tiny]{$-$} node[pos=.3,below,mklab]{$E$} (s1);
\end{tikzpicture}

\vspace{4mm}
\begin{columns}[T]
\begin{column}{.48\textwidth}
\footnotesize
\begin{itemize}\setlength\itemsep{1pt}
\item Two cascaded first-order lags: \hl{electrical} ($L,R$) and \hl{mechanical} ($J,B$).
\item The back-EMF $k_e$ is an \term{internal feedback loop} -- it gives the motor its natural speed stabilisation (damping).
\item $T_L$ enters as a \hl{disturbance} at the torque summing point.
\end{itemize}
\end{column}
\begin{column}{.48\textwidth}
\begin{notebox}
Servo drives close a fast \textbf{current loop} around $1/(Ls+R)$, so that $T_m\approx k_t\,i^*$ -- the motor then appears as a pure torque source acting on $1/(Js+B)$ (Day 4: cascaded current/speed/position loops).
\end{notebox}
\end{column}
\end{columns}
\end{frame}

% ---------------------------------------------------------------------
\begin{frame}{Transfer functions and time constants}
\begin{columns}[T]
\begin{column}{.55\textwidth}
From the block diagram ($T_L=0$):
\begin{eqbox}
$G_{\omega v}(s)=\dfrac{\Omega(s)}{V(s)}=\dfrac{k_t}{(Ls+R)(Js+B)+k_tk_e}$\\[2mm]
$\phantom{G_{\omega v}(s)}=\dfrac{k_t}{LJ\,s^2+(RJ+LB)\,s+(RB+k_tk_e)}$
\end{eqbox}
\vspace{1mm}
\footnotesize
Disturbance path:\quad $\dfrac{\Omega(s)}{T_L(s)}=-\dfrac{Ls+R}{(Ls+R)(Js+B)+k_tk_e}$\\[2mm]
Position:\quad $\dfrac{\Theta(s)}{V(s)}=\dfrac{1}{s}\,G_{\omega v}(s)$ \; (type-1, integrating)
\end{column}
\begin{column}{.43\textwidth}
\begin{defbox}[Time constants ($B\approx0$)]
\small
$\tau_e=\dfrac{L}{R}$ \quad electrical\\[1mm]
$\tau_m=\dfrac{R\,J}{k_t\,k_e}$ \quad mechanical
\end{defbox}
\footnotesize
\[
G_{\omega v}(s)\approx\frac{1/k_e}{\tau_m\tau_e s^2+\tau_m s+1}
\]
\[
\overset{\tau_m\gg\tau_e}{\approx}\;\frac{1/k_e}{(\tau_m s+1)(\tau_e s+1)}
\]
\begin{itemize}\setlength\itemsep{1pt}
\item Real poles if $\tau_m>4\tau_e$; otherwise oscillatory.
\item $\tau_m$ depends on \hl{total} inertia: load inertia slows the motor.
\item Typical: $\tau_e\approx$ 0.5--5\,ms, $\tau_m\approx$ 2--20\,ms.
\end{itemize}
\end{column}
\end{columns}
\end{frame}

% ---------------------------------------------------------------------
\begin{frame}{Worked example: from parameters to poles}
\begin{columns}[T]
\begin{column}{.40\textwidth}
\small
\renewcommand{\arraystretch}{1.1}
\begin{tabular}{@{}lr@{}}
\toprule
\textbf{Parameter} & \textbf{Value}\\
\midrule
Supply $v$ & \SI{48}{\volt}\\
$R$ & \SI{1.0}{\ohm}\\
$L$ & \SI{0.5}{\milli\henry}\\
$k_t=k_e$ & \SI{0.1}{\newton\metre\per\ampere}\\
$J$ (rotor) & \SI{1.0e-4}{\kilo\gram\metre\squared}\\
$B$ & \SI{5e-5}{\newton\metre\second\per\radian}\\
\bottomrule
\end{tabular}
\end{column}
\begin{column}{.58\textwidth}
\footnotesize
\begin{enumerate}\setlength\itemsep{2pt}
\item $\tau_e=L/R=\SI{0.5}{\milli\second}$, \quad $\tau_m=RJ/(k_tk_e)=\SI{10}{\milli\second}$ \;($\tau_m/\tau_e=20$).
\item Denominator coefficients:\\
$LJ=\num{5.0e-8}$, \; $RJ+LB=\num{1.0003e-4}$, \; $RB+k_tk_e=\num{1.005e-2}$.
\item DC gain: $k_t/(RB+k_tk_e)=\SI{9.95}{\radian\per\second\per\volt}$ ($\approx 1/k_e$).
\item Poles: $s^2+2000.5\,s+201\,000=0$\\
$\Rightarrow s_1=\SI{-106.1}{\per\second}$ ($\tau=\SI{9.42}{\milli\second}\approx\tau_m$),\\
$\phantom{\Rightarrow}\; s_2=\SI{-1894}{\per\second}$ ($\tau=\SI{0.53}{\milli\second}\approx\tau_e$).
\item No-load speed at \SI{48}{\volt}: $\omega_0=\SI{477.6}{\radian\per\second}=\SI{4561}{rpm}$.
\end{enumerate}
\begin{keybox}[Check]
Approximate $\tau_m,\tau_e$ are within 6\,\% of the exact pole time constants because $\tau_m\gg\tau_e$.
\end{keybox}
\end{column}
\end{columns}
\end{frame}

% ---------------------------------------------------------------------
\begin{frame}{Step response of the example motor (\SI{24}{\volt} step, $T_L=0$)}
\begin{columns}[T]
\begin{column}{.6\textwidth}
\begin{tikzpicture}
\begin{axis}[mkplot,width=8.2cm,height=5.4cm,xmin=0,xmax=50,ymin=0,ymax=260,
  xlabel={$t$ (ms)},ylabel={$\omega$ (\si{\radian\per\second})},axis y line*=left,axis x line*=bottom]
\addplot[mkBlue,thick,domain=0:50,samples=200]{238.81+14.169*exp(-1.8944*x)-252.97*exp(-0.10610*x)}; \label{pl:w}
\draw[dashed,mkGray] (axis cs:9.42,0) -- (axis cs:9.42,151) -- (axis cs:0,151);
\node[font=\tiny,text=mkGray,anchor=west] at (axis cs:9.8,140) {63\,\% at $\approx\tau_m$};
\end{axis}
\begin{axis}[mkplot,width=8.2cm,height=5.4cm,xmin=0,xmax=50,ymin=0,ymax=26,grid=none,
  axis y line*=right,axis x line=none,ylabel={$i$ (A)},legend pos=north east]
\addlegendimage{/pgfplots/refstyle=pl:w}\addlegendentry{speed $\omega$}
\addplot[mkOrange,thick,domain=0:50,samples=300]{0.1194-26.834*exp(-1.8944*x)+26.715*exp(-0.10610*x)}; \addlegendentry{current $i$}
\end{axis}
\end{tikzpicture}
\end{column}
\begin{column}{.38\textwidth}
\footnotesize
Exact solution (partial fractions):
\[
\omega(t)=238.8+14.2\,e^{-1894t}-253.0\,e^{-106.1t}
\]
\[
i(t)=0.12-26.8\,e^{-1894t}+26.7\,e^{-106.1t}
\]
\begin{itemize}\setlength\itemsep{1pt}
\item Current peaks at \hl{\SI{21.4}{\ampere}} after \SI{1.6}{\milli\second} -- 180$\times$ the no-load current!
\item Back-EMF builds up and pulls the current down: inherent damping.
\item Real drives \hl{limit the current}; with a current limit the speed ramps linearly ($\dot\omega=k_tI_\text{max}/J$).
\end{itemize}
\end{column}
\end{columns}
\end{frame}

% ---------------------------------------------------------------------
\begin{frame}{From datasheet to model -- units and pitfalls}
\footnotesize
\begin{columns}[T]
\begin{column}{.52\textwidth}
\renewcommand{\arraystretch}{1.15}
\begin{tabular}{@{}lll@{}}
\toprule
\textbf{Datasheet item} & \textbf{Typical unit} & \textbf{Model value}\\
\midrule
Back-EMF const. & V/krpm & $\div\,104.7\to\si{\volt\second\per\radian}$\\
Torque const. & \si{\newton\metre\per\ampere} (rms!) & $k_t$\\
Terminal resistance & \si{\ohm} line--line & $R_\text{ph}=R_{LL}/2$ (Y)\\
Terminal inductance & \si{\milli\henry} line--line & $L_\text{ph}=L_{LL}/2$ (Y)\\
Rotor inertia & \si{\kilo\gram\centi\metre\squared} & $\times10^{-4}\to\si{\kilo\gram\metre\squared}$\\
Mech. time const. & \si{\milli\second} & $\tau_m=RJ/(k_tk_e)$\\
Thermal resistance & \si{\kelvin\per\watt} & $R_\text{th}$ (Session 1.5)\\
\bottomrule
\end{tabular}
\end{column}
\begin{column}{.45\textwidth}
\begin{keybox}[Three-phase PMSM: watch the convention]
With $k_e$ in $\mathrm{V_{rms,LL}}/(\si{\radian\per\second})$ and $k_t$ in $\mathrm{N\,m/A_{rms}}$ (phase current):
\[
k_t=\sqrt3\,k_e
\]
$k_t=k_e$ holds only in \hl{consistent per-phase / dq units}.
\end{keybox}
\begin{itemize}\setlength\itemsep{1pt}
\item Example: \SI{33}{V_{rms}}/krpm $\Rightarrow$ $k_e=0.315$, $k_t=\SI{0.546}{\newton\metre\per\ampere}_\text{rms}$.
\item Hot winding: $R$ rises $\approx$ 0.39\,\%/K (copper): +39\,\% at $\Delta T=\SI{100}{\kelvin}$.
\item Datasheet values have \hl{$\pm$5--10\,\%} tolerance.
\end{itemize}
\end{column}
\end{columns}
\end{frame}

% ---------------------------------------------------------------------
\begin{frame}{Session 1.3 -- Summary / key take-aways}
\begin{columns}[T]
\begin{column}{.52\textwidth}
\begin{itemize}
\item Motor = RL circuit + back-EMF, coupled to an inertia with friction and load torque.
\item Energy conservation gives $k_t=k_e$ in SI units.
\item Transfer function $V\to\Omega$ is 2nd order with $\tau_e=L/R$ and $\tau_m=RJ/(k_tk_e)$.
\item Friction: viscous (linear) + Coulomb/stiction (nonlinear); $T_L$ is a disturbance.
\item Without current control, a voltage step causes large inrush current.
\end{itemize}
\end{column}
\begin{column}{.45\textwidth}
\begin{keybox}[Remember]
\small
$v=Ri+L\dot i+k_e\omega$\\[1mm]
$J\dot\omega=k_ti-B\omega-T_L$\\[1mm]
$\dfrac{\Omega}{V}=\dfrac{k_t}{(Ls+R)(Js+B)+k_tk_e}$\\[1mm]
$\tau_e=L/R,\quad \tau_m=RJ/(k_tk_e)$
\end{keybox}
\end{column}
\end{columns}
\end{frame}

% =====================================================================
\section{1.4 Actuators and motion transmission}
% =====================================================================

\begin{frame}{Motion transmission options}
\centering
\scriptsize
\renewcommand{\arraystretch}{1.05}
\begin{tabularx}{\textwidth}{@{}l>{\centering\arraybackslash}X>{\centering\arraybackslash}X>{\centering\arraybackslash}X>{\centering\arraybackslash}X>{\centering\arraybackslash}X@{}}
\toprule
\textbf{Mechanism} & \textbf{Conversion constant} & \textbf{Speed} & \textbf{Stiffness} & \textbf{Accuracy / backlash} & \textbf{Stroke}\\
\midrule
Ball screw & lead $p$ = 5--40\,mm/rev & $\le$ 1--2\,m/s & \textcolor{mkGreen}{high} & \textcolor{mkGreen}{$\mu$m, preloaded} & $\lesssim$ 3--4\,m (whip)\\
Toothed belt & pulley radius $r$ & \textcolor{mkGreen}{$\le$ 5--10\,m/s} & \textcolor{mkRed}{low--medium} & 0.05--0.1\,mm & long\\
Rack \& pinion & pinion radius $r$ & 2--5\,m/s & medium & backlash (dual pinion) & \textcolor{mkGreen}{unlimited}\\
Linear motor & direct, $k_f$ in N/A & \textcolor{mkGreen}{$>$5\,m/s} & servo-defined & \textcolor{mkGreen}{nm--$\mu$m (encoder)} & unlimited\\
Gearbox & ratio $N$ = 3--160 & motor-limited & medium & 1--10 arcmin & rotary\\
Direct-drive torque motor & 1:1 & low--medium & \textcolor{mkGreen}{highest} & \textcolor{mkGreen}{no backlash} & rotary\\
\bottomrule
\end{tabularx}

\vspace{1mm}
\begin{columns}[T]
\begin{column}{.48\textwidth}
\begin{defbox}[Geared / converted drive]
\footnotesize
$+$ small, fast motor; load inertia reduced by $1/N^2$\\
$-$ backlash, compliance, friction, wear, noise
\end{defbox}
\end{column}
\begin{column}{.48\textwidth}
\begin{defbox}[Direct drive]
\footnotesize
$+$ stiff, zero backlash, high bandwidth, low maintenance\\
$-$ large, expensive motor; full load inertia and disturbances act on the motor
\end{defbox}
\end{column}
\end{columns}
\end{frame}

% ---------------------------------------------------------------------
\begin{frame}{Linear electromechanical actuators: linear motor and voice coil}
\begin{columns}[T]
\begin{column}{.34\textwidth}
\centering
\figph{width=\linewidth,height=2.9cm,keepaspectratio}{d1_linear_motor.jpg}\\
\credit{Zureks, CC BY-SA 4.0, via Wikimedia Commons}
\vspace{1mm}
\footnotesize
\begin{itemize}\setlength\itemsep{1pt}
\item \term{Linear motor} = PMSM ``unrolled'': forcer (coil) + magnet track.
\item \term{Iron-core}: high force, cogging, attraction force $\approx$ 3--10$\times$ thrust.
\item \term{Ironless} (U-channel): no cogging, lower force, very smooth.
\end{itemize}
\end{column}
\begin{column}{.32\textwidth}
\centering
\begin{tikzpicture}[scale=.75]
\fill[mkGray!50] (0,0) rectangle (3.2,0.5);
\fill[mkGray!50] (0,0) rectangle (0.5,3.2);
\fill[mkGray!50] (1.2,0.5) rectangle (1.7,3.2);
\fill[mkRed!70] (0.5,0.5) rectangle (1.2,2.0); \node[white,font=\tiny\bfseries] at (0.85,1.25) {PM};
\fill[mkGray!50] (0.5,2.0) rectangle (1.2,3.2);
\fill[mkGray!50] (2.4,0.5) rectangle (3.2,3.2);
\fill[mkOrange] (1.85,1.4) rectangle (2.25,3.6);
\node[font=\tiny,text=mkOrange,rotate=90] at (2.05,2.5) {coil};
\foreach \y in {2.4,2.8} \draw[-{Stealth[length=1.3mm]},mkSky] (1.7,\y) -- (2.4,\y);
\node[font=\tiny,text=mkSky] at (2.6,3.75) {$B$};
\draw[-{Stealth},mkNavy,line width=1pt] (2.05,3.7) -- (2.05,4.4) node[above,font=\scriptsize]{$F=k_f\,i$};
\draw[dashed,mkGray] (0,-0.2) -- (0,4.6);
\node[font=\tiny,text=mkGray] at (1.6,-0.4) {axisymmetric half-section};
\end{tikzpicture}
\end{column}
\begin{column}{.32\textwidth}
\begin{eqbox}\small
$F=k_f\,i$,\qquad $e=k_f\,v$\\[1mm]
$k_f=N\,B\,l$\; [\si{\newton\per\ampere} = \si{\volt\second\per\metre}]\\[1mm]
$m\,\ddot x = k_f\, i - b\,\dot x - F_L$
\end{eqbox}
\footnotesize
\begin{itemize}\setlength\itemsep{1pt}
\item \term{Voice coil}: moving coil in a radial field -- no commutation, no cogging, stroke mm to a few cm.
\item Very high bandwidth: HDD head positioners, lens focus, fast tool servos, active vibration isolation.
\item Same equations as the DC motor with $T\to F$, $\omega\to v$, $J\to m$.
\end{itemize}
\end{column}
\end{columns}
\end{frame}

% ---------------------------------------------------------------------
\begin{frame}{Ball screw: kinematics and equivalent inertia}
\begin{columns}[T]
\begin{column}{.56\textwidth}
\centering
\begin{tikzpicture}[scale=.85]
\fill[mkBlue] (0,0.6) rectangle (1.4,2.0); \node[white,font=\scriptsize] at (0.7,1.3) {motor};
\node[font=\tiny] at (0.7,0.4) {$J_m$};
\fill[mkGray!70] (1.4,1.15) rectangle (1.9,1.45); \fill[mkAmber] (1.9,1.0) rectangle (2.3,1.6);
\node[font=\tiny] at (2.1,0.8) {$J_c$};
\fill[mkGray!40] (2.3,1.15) rectangle (8.2,1.45);
\foreach \x in {2.4,2.6,...,8.0} \draw[mkGray] (\x,1.15) -- (\x+0.12,1.45);
\fill[mkGray] (8.2,0.9) rectangle (8.5,1.7); \fill[mkGray] (2.3,0.9) rectangle (2.5,1.7);
\fill[mkOrange] (4.4,0.95) rectangle (5.4,1.65); \node[font=\tiny,white] at (4.9,1.3) {nut};
\fill[mkSky!70] (3.6,1.65) rectangle (6.2,2.3); \node[font=\scriptsize] at (4.9,1.98) {table $m$};
\draw[mkGray,line width=1pt] (2.2,0.2) -- (8.6,0.2);
\draw[-{Stealth},line width=1pt,mkNavy] (6.4,2.0) -- (7.3,2.0) node[right,font=\scriptsize]{$x,v$};
\draw[-{Stealth},line width=.8pt,mkOrange] (1.65,1.8) arc (100:-100:0.25 and 0.5);
\node[font=\scriptsize,text=mkOrange] at (1.6,2.45) {$\theta,\omega$};
\node[font=\tiny] at (5.6,0.6) {screw $J_s$, lead $p$};
\draw[-{Stealth},mkRed] (3.5,2.0) -- (2.8,2.0) node[left,font=\tiny,text=mkRed]{$F_L$};
\end{tikzpicture}

\vspace{1mm}
\figph{width=.6\linewidth,height=2.3cm,keepaspectratio}{d1_ballscrew.jpg}\\
\credit{Glenn McKechnie (Graibeard), CC BY-SA 3.0, via Wikimedia Commons}
\end{column}
\begin{column}{.42\textwidth}
\begin{eqbox}\small
$x=\dfrac{p}{2\pi}\,\theta$, \qquad $v=\dfrac{p}{2\pi}\,\omega$\\[1mm]
$T=\dfrac{p}{2\pi\,\eta}\,F$ \quad (driving)\\[1mm]
$J_\text{eq}=J_m+J_c+J_s+m\left(\dfrac{p}{2\pi}\right)^2$
\end{eqbox}
\footnotesize
\begin{itemize}\setlength\itemsep{1pt}
\item Derivation: $\tfrac12 m v^2=\tfrac12\,m\left(\tfrac{p}{2\pi}\right)^2\omega^2$ (kinetic energy).
\item Screw inertia: $J_s=\pi\rho L d^4/32$ -- grows with $d^4$!
\item $\eta\approx 0.9$ (ball screw) vs.\ 0.3--0.6 (sliding lead screw).
\item Limits: critical (whip) speed, buckling, nut DN value, axial stiffness.
\end{itemize}
\end{column}
\end{columns}
\end{frame}

% ---------------------------------------------------------------------
\begin{frame}{Gearboxes: planetary and harmonic (strain-wave) drives}
\begin{columns}[T]
\begin{column}{.25\textwidth}
\centering
\figph{width=\linewidth,height=2.5cm,keepaspectratio}{d1_planetary_diag.png}\\
\credit{Wapcaplet, CC BY-SA 3.0, via Wikimedia Commons}
\end{column}
\begin{column}{.25\textwidth}
\centering
\figph{width=\linewidth,height=2.5cm,keepaspectratio}{d1_harmonic_drive.jpg}\\
\credit{Pieceofmetalwork, CC BY-SA 4.0, via Wikimedia Commons}
\end{column}
\begin{column}{.47\textwidth}
\begin{eqbox}\footnotesize
\textbf{Planetary} (ring fixed, sun in, carrier out): $N=1+\dfrac{Z_\text{ring}}{Z_\text{sun}}$\\[1mm]
\textbf{Harmonic} (circular spline fixed, wave generator in, flexspline out):
$N=\dfrac{Z_\text{fs}}{Z_\text{cs}-Z_\text{fs}}$ \;(reversed)
\end{eqbox}
\footnotesize
Example: $Z_\text{fs}=200$, $Z_\text{cs}=202$ $\Rightarrow N=100$ in one stage.
\end{column}
\end{columns}
\vspace{1mm}
\scriptsize
\renewcommand{\arraystretch}{1.0}
\begin{tabularx}{\textwidth}{@{}l*{5}{>{\centering\arraybackslash}X}@{}}
\toprule
 & \textbf{Ratio / stage} & \textbf{Efficiency} & \textbf{Backlash} & \textbf{Torsional stiffness} & \textbf{Typical use}\\
\midrule
Planetary & 3--10 & \textcolor{mkGreen}{95--98\,\%} & 1--10 arcmin & high & servo axes, gantries\\
Harmonic drive & 30--160 & 60--90\,\% & \textcolor{mkGreen}{$\approx$ 0} & \textcolor{mkRed}{low, nonlinear} & robot joints\\
Cycloidal & 30--120 & 80--90\,\% & $<$ 1 arcmin & \textcolor{mkGreen}{very high} & heavy robot axes\\
Worm & 5--100 & \textcolor{mkRed}{40--90\,\%} & medium & medium & self-locking drives\\
\bottomrule
\end{tabularx}
\end{frame}

% ---------------------------------------------------------------------
\begin{frame}{Reflected inertia and load torque through a gear}
\begin{columns}[T]
\begin{column}{.46\textwidth}
\centering
\begin{tikzpicture}[scale=.9]
\fill[mkBlue] (0,0.4) rectangle (1.3,1.8); \node[white,font=\scriptsize] at (0.65,1.1) {$J_m$};
\draw[mkGray,line width=2pt] (1.3,1.1) -- (2.3,1.1);
\fill[mkSky] (2.3,0.75) rectangle (2.5,1.45); \node[font=\tiny] at (2.4,1.65) {$z_1$};
\fill[mkSky] (2.5,-1.1) rectangle (2.7,0.75); \node[font=\tiny] at (2.6,-1.3) {$z_2$};
\draw[mkGray,line width=2pt] (2.7,-0.2) -- (4.0,-0.2);
\fill[mkOrange] (4.0,-0.9) rectangle (4.6,0.5); \node[white,font=\scriptsize] at (4.3,-0.2) {$J_L$};
\node[font=\scriptsize,text=mkOrange] at (1.8,1.5) {$\omega_m,T_m$};
\node[font=\scriptsize,text=mkOrange] at (3.4,0.15) {$\omega_L$};
\draw[-{Stealth},mkRed] (5.2,0.3) arc (60:-60:0.5) node[right,font=\scriptsize,text=mkRed] {$T_L$};
\node[font=\scriptsize] at (2.4,-1.8) {$N=z_2/z_1=\omega_m/\omega_L$};
\end{tikzpicture}
\begin{tikzpicture}[scale=.9]
\fill[mkBlue] (0,0) rectangle (1.3,1.2); \node[white,font=\scriptsize] at (0.65,0.6) {$J_m$};
\draw[mkGray,line width=2pt] (1.3,0.6) -- (2.6,0.6);
\fill[mkOrange!60] (2.6,0.1) rectangle (3.6,1.1); \node[font=\scriptsize] at (3.1,0.6) {$J_L/N^2$};
\draw[-{Stealth},mkRed] (4.1,1.0) arc (60:-60:0.45) node[right,font=\scriptsize,text=mkRed] {$\dfrac{T_L}{N\eta}$};
\node[font=\scriptsize,text=mkGray] at (1.8,-0.3) {equivalent at motor shaft};
\end{tikzpicture}
\end{column}
\begin{column}{.52\textwidth}
\footnotesize
Kinetic energy must be preserved:
\[
\tfrac12 J_L\,\omega_L^2=\tfrac12\,\underbrace{\frac{J_L}{N^2}}_{J_L'}\,\omega_m^2
\]
Power balance with efficiency $\eta$:\quad $T_L'\,\omega_m\,\eta=T_L\,\omega_L\Rightarrow T_L'=\dfrac{T_L}{N\eta}$
\begin{eqbox}\small
$T_m=\Big(J_m+J_{g}+\dfrac{J_L}{N^2}\Big)\dot\omega_m+\dfrac{T_L}{N\eta}$
\end{eqbox}
\begin{tabular}{@{}ll@{}}
\toprule
\textbf{Mechanism} & \textbf{Reflected inertia at motor}\\
\midrule
Gear ratio $N$ & $J/N^2$\\
Ball screw, lead $p$ & $m\,(p/2\pi)^2$\\
Belt / rack, radius $r$ & $m\,r^2$ (+ pulleys/pinion)\\
\bottomrule
\end{tabular}
\end{column}
\end{columns}
\end{frame}

% ---------------------------------------------------------------------
\begin{frame}{Inertia ratio and the optimal gear ratio}
\begin{columns}[T]
\begin{column}{.52\textwidth}
\begin{tikzpicture}
\begin{axis}[mkplot,width=7.4cm,height=5.2cm,xmin=0.5,xmax=20,ymin=0,ymax=1.08,xmode=log,
  xlabel={gear ratio $N$},ylabel={$\alpha_L/\alpha_{L,\max}$},legend pos=south east,
  xtick={1,2,4,8,16},xticklabels={1,2,4,8,16}]
\addplot[mkSky,thick,domain=0.5:20,samples=120]{2*x*2/(x^2+4)}; \addlegendentry{$J_L/J_m=4$}
\addplot[mkBlue,thick,domain=0.5:20,samples=120]{2*x*4/(x^2+16)}; \addlegendentry{$J_L/J_m=16$}
\addplot[mkOrange,thick,domain=0.5:20,samples=120]{2*x*8/(x^2+64)}; \addlegendentry{$J_L/J_m=64$}
\draw[dashed,mkGray] (axis cs:2,0) -- (axis cs:2,1);
\draw[dashed,mkGray] (axis cs:4,0) -- (axis cs:4,1);
\draw[dashed,mkGray] (axis cs:8,0) -- (axis cs:8,1);
\end{axis}
\end{tikzpicture}
\end{column}
\begin{column}{.46\textwidth}
\footnotesize
Load acceleration for a given motor torque $T_m$ (no friction):
\[
\alpha_L=\frac{N\,T_m}{J_mN^2+J_L}
\]
$d\alpha_L/dN=0$ gives
\begin{eqbox}\small
$N_\text{opt}=\sqrt{J_L/J_m}$ \;$\Leftrightarrow$\; $J_L/N^2=J_m$ (\hl{inertia matching})
\end{eqbox}
\begin{itemize}\setlength\itemsep{1pt}
\item The maximum is flat: $N$ within $\pm$30\,\% of $N_\text{opt}$ costs $<$\,6\,\%.
\item In practice $N$ is set by \hl{speed}: $N\le\omega_{m,\max}/\omega_{L,\max}$.
\item \term{Inertia ratio} $J_L'/J_m$: guideline $\le$ 3--10 for good servo tuning with compliant mechanics; direct drives tolerate more.
\item Large ratios $\Rightarrow$ low mechanical resonance relative to bandwidth, sensitive tuning (Day 5).
\end{itemize}
\end{column}
\end{columns}
\end{frame}

% ---------------------------------------------------------------------
\begin{frame}{Transmission non-idealities: efficiency, backlash, stiffness}
\begin{columns}[T]
\begin{column}{.33\textwidth}
\centering
\begin{tikzpicture}
\begin{axis}[mkplot,width=5.0cm,height=4.4cm,xmin=-2,xmax=2,ymin=-1.6,ymax=1.6,axis lines=middle,
  xlabel={$\theta_m/N$},ylabel={$\theta_L$},xtick=\empty,ytick=\empty,
  every axis x label/.style={at={(current axis.right of origin)},anchor=north east,font=\scriptsize},
  every axis y label/.style={at={(current axis.above origin)},anchor=north west,font=\scriptsize}]
\addplot[mkGray,thin,dashed,domain=-1.2:1.9]{x-0.4};
\addplot[mkGray,thin,dashed,domain=-1.9:1.2]{x+0.4};
\addplot[mkOrange,thick,-{Stealth}] coordinates {(-1.0,-1.4) (1.4,1.0)};
\addplot[mkOrange,thick,-{Stealth}] coordinates {(1.4,1.0) (0.6,1.0)};
\addplot[mkOrange,thick,-{Stealth}] coordinates {(0.6,1.0) (-1.0,-0.6)};
\draw[mkNavy,{Stealth[length=1mm]}-{Stealth[length=1mm]}] (axis cs:0.6,1.2) -- (axis cs:1.4,1.2);
\node[font=\tiny,text=mkNavy] at (axis cs:1.0,1.42) {$2\Delta$};
\end{axis}
\end{tikzpicture}\\
{\tiny backlash: dead zone of width $2\Delta$}
\end{column}
\begin{column}{.33\textwidth}
\centering
\begin{tikzpicture}[scale=.8]
\fill[mkBlue] (0,0) rectangle (1,1.2); \node[white,font=\scriptsize] at (0.5,0.6) {$J_m$};
\draw[decorate,decoration={coil,aspect=0.5,segment length=2mm,amplitude=2mm},mkNavy] (1,0.6) -- (2.6,0.6);
\node[font=\scriptsize] at (1.8,1.1) {$k$};
\fill[mkOrange] (2.6,0) rectangle (3.6,1.2); \node[white,font=\scriptsize] at (3.1,0.6) {$J_L'$};
\end{tikzpicture}
\[
f_r=\frac{1}{2\pi}\sqrt{\frac{k\,(J_m+J_L')}{J_m\,J_L'}}
\]
{\tiny compliant gear/coupling/belt $\Rightarrow$ resonance (Session 2.1, Day 5)}
\end{column}
\begin{column}{.32\textwidth}
\footnotesize
\begin{itemize}\setlength\itemsep{2pt}
\item \term{Efficiency} $\eta$: losses $\propto$ torque; back-driving efficiency differs (worm/lead screw may self-lock).
\item \term{Backlash} $\Delta$: lost motion at reversal, limit cycles in position loops; reduce by preload, dual pinions, harmonic drives.
\item \term{Stiffness} $k$: belts and harmonic drives are compliant; ball screws stiff axially but screw torsion/nut stiffness matter.
\item Load-side encoder ``sees'' backlash and compliance -- dual-loop control (Day 5).
\end{itemize}
\end{column}
\end{columns}
\end{frame}

% ---------------------------------------------------------------------
\begin{frame}{Session 1.4 -- Summary / key take-aways}
\begin{columns}[T]
\begin{column}{.54\textwidth}
\begin{itemize}
\item Transmissions trade motor size against stiffness, backlash and friction.
\item Every load is reflected to the motor shaft by \hl{energy equivalence}.
\item Optimal (acceleration) gear ratio: $N_\text{opt}=\sqrt{J_L/J_m}$; in practice speed limits $N$.
\item Keep the reflected inertia ratio moderate; compliant drives need careful tuning.
\item Linear motors and voice coils obey the same equations with $F,\,v,\,m$.
\end{itemize}
\end{column}
\begin{column}{.42\textwidth}
\begin{keybox}[Remember]
\small
$J_L'=J_L/N^2$, \; $T_L'=T_L/(N\eta)$\\[1mm]
$J_\text{screw}'=m\,(p/2\pi)^2$, \; $J_\text{belt}'=m\,r^2$\\[1mm]
$N_\text{opt}=\sqrt{J_L/J_m}$\\[1mm]
$f_r=\frac{1}{2\pi}\sqrt{k(J_m+J_L')/(J_mJ_L')}$
\end{keybox}
\end{column}
\end{columns}
\end{frame}

% =====================================================================
\section{1.5 Characteristics and selection criteria}
% =====================================================================

\begin{frame}{Servo motor torque--speed envelope}
\begin{columns}[T]
\begin{column}{.58\textwidth}
\begin{tikzpicture}
\begin{axis}[mkplot,width=8cm,height=5.6cm,xmin=0,xmax=5200,ymin=0,ymax=4.3,
  xlabel={speed $n$ (rpm)},ylabel={torque $T$ (\si{\newton\metre})},scaled x ticks=false,
  xtick={0,1000,2000,3000,4000,5000}]
\addplot[fill=mkCream,draw=none] coordinates {(0,1.27)(0,3.82)(3300,3.82)(5000,1.6)(5000,0.95)(3000,1.27)} -- cycle;
\addplot[fill=mkPale,draw=none] coordinates {(0,0)(0,1.27)(3000,1.27)(5000,0.95)(5000,0)} -- cycle;
\addplot[mkOrange,thick] coordinates {(0,3.82)(3300,3.82)(5000,1.6)};
\addplot[mkBlue,thick] coordinates {(0,1.27)(3000,1.27)(5000,0.95)};
\node[font=\scriptsize,text=mkBlue] at (axis cs:1500,0.7) {continuous (S1) zone};
\node[font=\scriptsize,text=mkOrange] at (axis cs:1500,2.6) {intermittent zone};
\node[font=\tiny,text=mkOrange,anchor=south west] at (axis cs:4250,2.7) {voltage limit};
\node[font=\tiny,text=mkNavy,anchor=south] at (axis cs:1000,3.85) {$T_\text{peak}$ (current limit)};
\node[font=\tiny,text=mkNavy,anchor=south] at (axis cs:4000,1.3) {$T_\text{cont}$};
\node[circle,fill=mkNavy,inner sep=1.2pt] at (axis cs:3000,1.27) {};
\node[font=\tiny,anchor=north] at (axis cs:3000,1.2) {rated point};
\end{axis}
\end{tikzpicture}\\
{\tiny\color{mkGray} Generic 400\,W AC servo class: $T_N=\SI{1.27}{\newton\metre}$, $T_\text{peak}=\SI{3.82}{\newton\metre}$, $n_N=\SI{3000}{rpm}$.}
\end{column}
\begin{column}{.40\textwidth}
\footnotesize
\begin{itemize}\setlength\itemsep{2pt}
\item \term{Continuous torque} $T_\text{cont}$: limited by \hl{winding temperature} at steady state (S1); decreases at speed (iron and friction losses).
\item \term{Peak torque} $T_\text{peak}$: limited by drive current, magnet demagnetisation and saturation; allowed for short times only.
\item \term{Voltage limit}: at high speed the back-EMF leaves no voltage reserve to force current $\Rightarrow$ $T_\text{peak}$ drops. Depends on \hl{DC bus voltage}.
\item Curves are specified for a given ambient (40\,$^\circ$C), mounting flange and drive.
\end{itemize}
\end{column}
\end{columns}
\end{frame}

% ---------------------------------------------------------------------
\begin{frame}{Thermal limits: thermal model and time constant}
\begin{columns}[T]
\begin{column}{.40\textwidth}
\centering
\begin{circuitikz}[scale=.8,transform shape,font=\small]
\draw (0,0) to[I,l=$P_\text{loss}$,color=mkOrange] (0,2.5) -- (2,2.5)
 to[R,l=$R_\text{th}$,color=mkNavy] (2,0) -- (0,0);
\draw (2,2.5) -- (4,2.5) to[C,l=$C_\text{th}$,color=mkNavy] (4,0) -- (2,0);
\node[font=\scriptsize,above] at (3,2.5) {$\Delta\vartheta$};
\draw (0,0) node[ground]{};
\node[font=\scriptsize,text=mkGray] at (2.2,-0.8) {ambient = reference};
\end{circuitikz}
\begin{eqbox}\footnotesize
$C_\text{th}\,\dfrac{d\Delta\vartheta}{dt}=P_\text{loss}-\dfrac{\Delta\vartheta}{R_\text{th}}$\\[1mm]
$\Delta\vartheta(t)=P\,R_\text{th}\big(1-e^{-t/\tau_\text{th}}\big)$\\[1mm]
$\tau_\text{th}=R_\text{th}\,C_\text{th}$
\end{eqbox}
\end{column}
\begin{column}{.58\textwidth}
\begin{tikzpicture}
\begin{axis}[mkplot,width=7.4cm,height=4.4cm,xmin=0,xmax=3,ymin=0,ymax=1.6,
  xlabel={$t/\tau_\text{th}$},ylabel={$\Delta\vartheta/\Delta\vartheta_\text{max}$},legend pos=south east]
\addplot[mkBlue,thick,domain=0:3,samples=60]{1-exp(-x)}; \addlegendentry{$T=T_N$}
\addplot[mkSky,thick,domain=0:0.75,samples=40]{2.25*(1-exp(-x))}; \addlegendentry{$T=1.5\,T_N$}
\addplot[mkOrange,thick,domain=0:0.5,samples=40]{4*(1-exp(-x))}; \addlegendentry{$T=2\,T_N$}
\addplot[mkRed,dashed,domain=0:3]{1};
\draw[mkGray,dotted] (axis cs:0.288,0) -- (axis cs:0.288,1);
\draw[mkGray,dotted] (axis cs:0.588,0) -- (axis cs:0.588,1);
\node[font=\tiny,anchor=south] at (axis cs:0.22,1.02) {0.29};
\node[font=\tiny,anchor=south] at (axis cs:0.66,1.02) {0.59};
\end{axis}
\end{tikzpicture}
\footnotesize
\begin{itemize}\setlength\itemsep{1pt}
\item Copper loss $\propto T^2$: overload $k=T/T_N$ reaches the limit after $t=-\tau_\text{th}\ln(1-1/k^2)$.
\item $\tau_\text{th}$ (motor) $\approx$ 10--60 min, but the \hl{winding} has seconds-scale time constants: peak torque only for $\approx$ 0.1--few s.
\item Insulation class F/H (155/180\,$^\circ$C); magnets degrade if too hot. Drives use $I^2t$ models and sensors (PTC/KTY/Pt1000).
\end{itemize}
\end{column}
\end{columns}
\end{frame}

% ---------------------------------------------------------------------
\begin{frame}{Duty cycle and RMS torque}
\begin{columns}[T]
\begin{column}{.52\textwidth}
\begin{tikzpicture}
\begin{groupplot}[group style={group size=1 by 2,vertical sep=4mm,xticklabels at=edge bottom},
  mkplot,width=7.2cm,height=3.2cm,xmin=0,xmax=1.0]
\nextgroupplot[ymin=0,ymax=3500,ylabel={$n$ (rpm)},scaled y ticks=false,ytick={0,1500,3000}]
\addplot[mkBlue,thick] coordinates {(0,0)(0.1,3000)(0.4,3000)(0.5,0)(1.0,0)};
\nextgroupplot[ymin=-1.1,ymax=1.2,ylabel={$T$ (\si{\newton\metre})},xlabel={$t$ (s)}]
\addplot[mkOrange,thick,const plot] coordinates {(0,0.924)(0.1,0.059)(0.4,-0.806)(0.5,0)(1.0,0)};
\addplot[mkSky,thick,dashed,domain=0:1]{0.389};
\addplot[mkSky,thick,dashed,domain=0:1]{-0.389};
\node[font=\tiny,text=mkMid,anchor=west] at (axis cs:0.62,0.55) {$T_\text{rms}=0.389$};
\node[font=\tiny] at (axis cs:0.05,0.6) {$T_1$};
\node[font=\tiny] at (axis cs:0.25,0.25) {$T_2$};
\node[font=\tiny] at (axis cs:0.45,-0.55) {$T_3$};
\node[font=\tiny] at (axis cs:0.75,-0.2) {dwell};
\end{groupplot}
\end{tikzpicture}
\end{column}
\begin{column}{.46\textwidth}
\footnotesize
Copper losses $\propto i^2\propto T^2$. If the cycle time $t_\text{cyc}\ll\tau_\text{th}$, the winding sees the \hl{mean loss}:
\begin{eqbox}\small
$T_\text{rms}=\sqrt{\dfrac{1}{t_\text{cyc}}\displaystyle\sum_k T_k^2\,t_k}\;\le\;T_\text{cont}(\bar n)$
\end{eqbox}
\begin{itemize}\setlength\itemsep{1pt}
\item Trapezoidal profile: $T_1=J_\text{tot}\alpha+T_f$, $T_2=T_f$, $T_3=-J_\text{tot}\alpha+T_f$, dwell: $T_4=T_\text{hold}$.
\item Include holding torque on \hl{vertical axes} (gravity) -- dwell is not free!
\item Duty types per IEC 60034-1: S1 continuous, S2 short-time, S3 intermittent periodic, \dots
\item Also check \hl{peak} $\max|T_k|\le T_\text{peak}(n)$ and $n_\text{max}$.
\end{itemize}
\end{column}
\end{columns}
\end{frame}

% ---------------------------------------------------------------------
\begin{frame}{Worked sizing example (1/2): load and move profile}
\begin{columns}[T]
\begin{column}{.40\textwidth}
\begin{defbox}[Given: horizontal ball-screw axis]
\footnotesize
Table + workpiece $m=\SI{50}{\kilo\gram}$\\
Screw: lead $p=\SI{10}{\milli\metre}$, $d=\SI{20}{\milli\metre}$, $L=\SI{0.8}{\metre}$, steel, $\eta=0.9$\\
Coupling $J_c=\SI{0.10e-4}{\kilo\gram\metre\squared}$\\
Guide friction $\mu=0.01$; bearings/seal $T_b=\SI{0.05}{\newton\metre}$\\
Move \SI{200}{\milli\metre}: $t_a=t_d=\SI{0.1}{\second}$, total \SI{0.5}{\second}, dwell \SI{0.5}{\second} ($t_\text{cyc}=\SI{1}{\second}$)
\end{defbox}
\end{column}
\begin{column}{.58\textwidth}
\footnotesize
\textbf{1. Kinematics}\\
$v_\text{max}=\dfrac{s}{t_a+t_c}=\dfrac{0.2}{0.4}=\SI{0.5}{\metre\per\second}$, \;
$n=\dfrac{v}{p}=\SI{50}{rev/s}=\SI{3000}{rpm}$\\
$\omega_\text{max}=\SI{314.2}{\radian\per\second}$, \; $\alpha=\omega_\text{max}/t_a=\SI{3142}{\radian\per\second\squared}$\\[1.5mm]
\textbf{2. Inertia at the motor shaft}\\
$J_s=\pi\rho L d^4/32=\pi\cdot7850\cdot0.8\cdot0.02^4/32=\SI{0.99e-4}{\kilo\gram\metre\squared}$\\
$J_t=m\,(p/2\pi)^2=50\,(0.01/2\pi)^2=\SI{1.27e-4}{\kilo\gram\metre\squared}$\\
$J_L=J_s+J_t+J_c=\SI{2.35e-4}{\kilo\gram\metre\squared}$\\[1.5mm]
\textbf{3. Friction torque}\\
$T_f=\dfrac{\mu m g\,p}{2\pi\eta}+T_b=\dfrac{0.01\cdot50\cdot9.81\cdot0.01}{2\pi\cdot0.9}+0.05=\SI{0.059}{\newton\metre}$
\begin{notebox}
Assumption: efficiency applied to friction only, not to inertia (conservative alternative: $J_t/\eta$).
\end{notebox}
\end{column}
\end{columns}
\end{frame}

% ---------------------------------------------------------------------
\begin{frame}{Worked sizing example (2/2): torques, RMS and check}
\begin{columns}[T]
\begin{column}{.52\textwidth}
\footnotesize
\textbf{4. Candidate:} 400\,W servo, $J_m=\SI{0.40e-4}{\kilo\gram\metre\squared}$, $T_N=1.27$, $T_\text{peak}=\SI{3.82}{\newton\metre}$, $n_N=\SI{3000}{rpm}$\\
Inertia ratio $J_L/J_m=5.9$ \;\textcolor{mkGreen}{\checkmark}\quad $J_\text{tot}=\SI{2.75e-4}{\kilo\gram\metre\squared}$\\[1.5mm]
\textbf{5. Segment torques}\\
$J_\text{tot}\alpha=2.75\cdot10^{-4}\cdot3142=\SI{0.865}{\newton\metre}$\\
$T_1=0.865+0.059=\SI{0.924}{\newton\metre}$ \;(0.1\,s)\\
$T_2=\SI{0.059}{\newton\metre}$ \;(0.3\,s), \quad $T_3=-0.865+0.059=\SI{-0.806}{\newton\metre}$ \;(0.1\,s), \quad $T_4=0$\\[1.5mm]
\textbf{6. RMS torque}\\
$T_\text{rms}=\sqrt{\dfrac{0.924^2\cdot0.1+0.059^2\cdot0.3+0.806^2\cdot0.1}{1.0}}$ $=$ \hl{\SI{0.389}{\newton\metre}}\\[1.5mm]
\textbf{7. Checks}: peak $0.924<3.82$ \textcolor{mkGreen}{\checkmark}; \; RMS $0.389<1.27$ (31\,\%) \textcolor{mkGreen}{\checkmark}; \; $n_\text{max}=3000\le n_N$ \textcolor{mkGreen}{\checkmark}\\
Regenerated energy per stop: $\tfrac12J_\text{tot}\omega^2=\SI{13.6}{\joule}$ (bus capacitor / brake resistor).
\end{column}
\begin{column}{.46\textwidth}
\begin{tikzpicture}
\begin{axis}[mkplot,width=6.6cm,height=4.6cm,xmin=0,xmax=5200,ymin=0,ymax=4.3,
  xlabel={$n$ (rpm)},ylabel={$|T|$ (\si{\newton\metre})},scaled x ticks=false,xtick={0,2000,4000}]
\addplot[fill=mkCream,draw=none] coordinates {(0,1.27)(0,3.82)(3300,3.82)(5000,1.6)(5000,0.95)(3000,1.27)} -- cycle;
\addplot[fill=mkPale,draw=none] coordinates {(0,0)(0,1.27)(3000,1.27)(5000,0.95)(5000,0)} -- cycle;
\addplot[mkOrange,thick] coordinates {(0,3.82)(3300,3.82)(5000,1.6)};
\addplot[mkBlue,thick] coordinates {(0,1.27)(3000,1.27)(5000,0.95)};
\addplot[mkRed,thick,mark=*,mark size=1.2pt] coordinates {(0,0.924)(3000,0.924)};
\addplot[mkMid,thick,mark=*,mark size=1.2pt] coordinates {(3000,0.806)(0,0.806)};
\addplot[only marks,mark=diamond*,mark size=2.5pt,mkNavy] coordinates {(1200,0.389)};
\node[font=\tiny,text=mkRed,anchor=south] at (axis cs:1500,0.93) {accel.};
\node[font=\tiny,text=mkNavy,anchor=west] at (axis cs:1300,0.33) {$T_\text{rms}$ at $\bar n=1200$};
\end{axis}
\end{tikzpicture}
\begin{notebox}
A 200\,W motor ($J_m=\SI{0.18e-4}{}$, $T_N=0.64$, $T_\text{pk}=\SI{1.91}{\newton\metre}$) gives $T_\text{rms}=0.358$, $T_\text{pk}=0.854$ -- thermally OK, but inertia ratio 13: \hl{tuning/stiffness decides}.
\end{notebox}
\end{column}
\end{columns}
\end{frame}

% ---------------------------------------------------------------------
\begin{frame}{Motor--drive matching and selection checklist}
\begin{columns}[T]
\begin{column}{.5\textwidth}
\footnotesize
\textbf{Voltage}: the drive must supply back-EMF + drops at top speed
\begin{eqbox}\small
$V_\text{req}\approx k_e\,\omega_\text{max}+R\,I_\text{pk}\;(+\,\omega_eL\,I)\;\le\;V_\text{avail}(V_\text{DC})$
\end{eqbox}
\textbf{Current}: $I_\text{pk}=T_\text{pk}/k_t$, \; $I_\text{rms}=T_\text{rms}/k_t$ $\le$ drive peak / continuous ratings.\\[1mm]
\textbf{Example} (motor above, $k_e=\SI{33}{V_{rms,LL}}$/krpm, $k_t=\SI{0.546}{\newton\metre\per\ampere}$):
\begin{itemize}\setlength\itemsep{1pt}
\item EMF at 3000 rpm: \SI{99}{V_{rms,LL}} $\Rightarrow V_\text{DC}\gtrsim\sqrt2\cdot99\approx\SI{140}{\volt}$ + margin.
\item 230\,V drive ($V_\text{DC}\approx\SI{325}{\volt}$): \textcolor{mkGreen}{\checkmark}; \; 48\,V drive: $n_\text{max}\approx\SI{1000}{rpm}$ \textcolor{mkRed}{$\times$} $\Rightarrow$ low-voltage winding needed.
\item $I_\text{pk}=0.924/0.546=\SI{1.7}{A_{rms}}$, \; $I_\text{rms}=\SI{0.71}{A_{rms}}$.
\end{itemize}
\textbf{Same motor frame, different winding}: $k_t\uparrow$ $\Leftrightarrow$ less current, but higher voltage per speed.
\end{column}
\begin{column}{.47\textwidth}
\begin{keybox}[Selection checklist]
\footnotesize
\begin{enumerate}\setlength\itemsep{0pt}
\item Define move profile, loads, cycle, environment.
\item Reflect inertia/torques to the motor (gear, screw, belt).
\item Choose ratio: speed range, $N_\text{opt}$, resolution.
\item Peak torque at all speeds inside the \hl{intermittent zone}.
\item $T_\text{rms}$ at mean speed inside the \hl{continuous zone} (with 20--30\,\% margin).
\item Inertia ratio acceptable for the mechanics' stiffness.
\item Drive: voltage, peak/continuous current, regen energy.
\item Feedback resolution, brake, IP class, standards (IEC 60034, IEC 61800).
\end{enumerate}
\end{keybox}
\end{column}
\end{columns}
\end{frame}

% ---------------------------------------------------------------------
\begin{frame}{Session 1.5 -- Summary / key take-aways}
\begin{columns}[T]
\begin{column}{.54\textwidth}
\begin{itemize}
\item The torque--speed envelope has a \hl{thermal} (continuous) and a \hl{current/voltage} (intermittent) boundary.
\item Heat accumulates with $\tau_\text{th}$ of minutes: judge with RMS torque over the cycle.
\item Peak torque must be available at the speed where it is needed.
\item Motor and drive are matched through $k_t$, $k_e$, bus voltage and current ratings.
\item Size systematically: kinematics $\to$ inertias $\to$ torques $\to$ RMS/peak $\to$ drive.
\end{itemize}
\end{column}
\begin{column}{.42\textwidth}
\begin{keybox}[Remember]
\small
$T_\text{rms}=\sqrt{\tfrac{1}{t_\text{cyc}}\sum T_k^2t_k}$\\[1mm]
$\Delta\vartheta=P R_\text{th}(1-e^{-t/\tau_\text{th}})$\\[1mm]
$I_\text{pk}=T_\text{pk}/k_t$, \; $V\gtrsim k_e\omega_\text{max}+RI_\text{pk}$
\end{keybox}
\end{column}
\end{columns}
\end{frame}

% =====================================================================
\section{2.1 Modelling of dynamic systems -- fundamental physical models}
% =====================================================================

\begin{frame}{Lumped-parameter modelling: workflow and building blocks}
\begin{columns}[T]
\begin{column}{.40\textwidth}
\centering
\begin{tikzpicture}[node distance=2.2mm,
  st/.style={mkblk,minimum width=40mm,text width=38mm,minimum height=6mm}]
\node[st] (a) {1. System boundary, inputs, outputs};
\node[st,below=of a] (b) {2. Idealise into lumped elements};
\node[st,below=of b] (c) {3. Free-body diagrams / Kirchhoff laws};
\node[st,below=of c] (d) {4. Constitutive laws of elements};
\node[st,below=of d,fill=mkCream,draw=mkOrange] (e) {5. ODEs $\to$ TF / state space};
\node[st,below=of e] (f) {6. Validate against measurement};
\foreach \x/\y in {a/b,b/c,c/d,d/e,e/f} \draw[mksig] (\x) -- (\y);
\draw[mksig,mkGray,dashed] (f.east) -- ++(4mm,0) |- (b.east);
\end{tikzpicture}
\end{column}
\begin{column}{.58\textwidth}
\footnotesize
\renewcommand{\arraystretch}{1.15}
\begin{tabular}{@{}llll@{}}
\toprule
\textbf{Domain} & \textbf{Kinetic store} & \textbf{Potential store} & \textbf{Dissipator}\\
\midrule
Translation & mass $m$ & spring $k$ & damper $b$\\
Rotation & inertia $J$ & torsion spring $k_\theta$ & damper $B$\\
Electrical & inductor $L$ & capacitor $C$ & resistor $R$\\
\bottomrule
\end{tabular}

\vspace{2mm}
\begin{itemize}\setlength\itemsep{1pt}
\item \term{Lumped}: wave propagation neglected -- valid for frequencies well below the first ignored structural mode.
\item Each independent energy store adds one state ($\Rightarrow$ system order).
\item \term{Transformers} (gears, screws) and \term{gyrators} (motors: $T=k i$, $e=k\omega$) connect domains without storing energy.
\item Choose the \hl{simplest model that answers the question} (Day 2: linearisation, identification).
\end{itemize}
{\scriptsize\color{mkGray}[Karnopp, Margolis \& Rosenberg, 2012; Ogata, 2010]}
\end{column}
\end{columns}
\end{frame}

% ---------------------------------------------------------------------
\begin{frame}{Mass--spring--damper: free-body diagram to transfer function}
\begin{columns}[T]
\begin{column}{.48\textwidth}
\centering
\begin{tikzpicture}[scale=.9]
\fill[pattern=north east lines,pattern color=mkGray] (-0.3,-0.2) rectangle (0,2.2); \draw[mkGray,line width=1pt] (0,-0.2) -- (0,2.2);
\draw[decorate,decoration={coil,aspect=0.5,segment length=2.4mm,amplitude=2.2mm},mkNavy,line width=.6pt] (0,1.6) -- (2,1.6);
\node[font=\scriptsize] at (1,2.05) {$k$};
\draw[mkNavy,line width=.6pt] (0,0.4) -- (0.9,0.4); \draw[mkNavy,line width=.6pt] (0.9,0.2) -- (0.9,0.6);
\draw[mkNavy,line width=.6pt] (0.7,0.15) -- (1.3,0.15) -- (1.3,0.65) -- (0.7,0.65); \draw[mkNavy,line width=.6pt] (1.3,0.4) -- (2,0.4);
\node[font=\scriptsize] at (1,-0.1) {$b$};
\fill[mkSky!60,draw=mkBlue] (2,-0.1) rectangle (3.2,2.1); \node[font=\small] at (2.6,1) {$m$};
\draw[-{Stealth},mkOrange,line width=1pt] (3.2,1) -- (4.1,1) node[right,font=\scriptsize]{$F(t)$};
\draw[-{Stealth},mkNavy] (2,2.4) -- (2.9,2.4) node[right,font=\scriptsize]{$x$};
\draw[mkNavy] (2,2.25) -- (2,2.55);
\fill[pattern=north east lines,pattern color=mkGray] (1.8,-0.45) rectangle (3.4,-0.2);
\begin{scope}[shift={(0.8,-2.2)}]
\fill[mkSky!60,draw=mkBlue] (1.2,0) rectangle (2.4,1.2); \node[font=\small] at (1.8,0.6) {$m$};
\draw[-{Stealth},mkRed,line width=.8pt] (1.2,0.85) -- (0.1,0.85) node[left,font=\scriptsize]{$kx$};
\draw[-{Stealth},mkRed,line width=.8pt] (1.2,0.35) -- (0.1,0.35) node[left,font=\scriptsize]{$b\dot x$};
\draw[-{Stealth},mkOrange,line width=.8pt] (2.4,0.6) -- (3.4,0.6) node[right,font=\scriptsize]{$F$};
\node[font=\scriptsize,text=mkGray] at (1.8,-0.3) {free-body diagram};
\end{scope}
\end{tikzpicture}
\end{column}
\begin{column}{.50\textwidth}
\footnotesize
Newton: $\sum F = m\ddot x$ \; (positive $x$ to the right):
\begin{eqbox}
$m\,\ddot x + b\,\dot x + k\,x = F(t)$\\[1mm]
$G(s)=\dfrac{X(s)}{F(s)}=\dfrac{1}{ms^2+bs+k}$\\[1mm]
$\phantom{G(s)}=\dfrac{1}{k}\,\dfrac{\omega_n^2}{s^2+2\zeta\omega_n s+\omega_n^2}$
\end{eqbox}
\begin{defbox}[Characteristic parameters]
\small
$\omega_n=\sqrt{k/m}$ \quad natural frequency\\
$\zeta=\dfrac{b}{2\sqrt{km}}$ \quad damping ratio\\
$\omega_d=\omega_n\sqrt{1-\zeta^2}$ \quad damped frequency
\end{defbox}
Example: $m=\SI{2}{\kilo\gram}$, $k=\SI{2e5}{\newton\per\metre}$, $b=\SI{120}{\newton\second\per\metre}$ $\Rightarrow$ $\omega_n=\SI{316}{\radian\per\second}$ (\SI{50.3}{\hertz}), $\zeta=0.095$.
\end{column}
\end{columns}
\end{frame}

% ---------------------------------------------------------------------
\begin{frame}{Second-order step responses and pole locations}
\begin{columns}[T]
\begin{column}{.60\textwidth}
\begin{tikzpicture}
\begin{axis}[mkplot,width=8.4cm,height=5.6cm,xmin=0,xmax=14,ymin=0,ymax=1.85,
  xlabel={normalised time $\omega_n t$},ylabel={$x(t)\,k/F_0$},legend pos=north east,legend columns=2,
  cycle list={{mkOrange,thick},{mkAmber,thick},{mkBlue,thick},{mkMid,thick,dashed},{mkSky,thick}}]
\addplot+[mark=none,domain=0:14,samples=300]{1-exp(-0.1*x)/sqrt(1-0.01)*sin(deg(sqrt(1-0.01)*x)+deg(acos(0.1)))}; \addlegendentry{$\zeta=0.1$}
\addplot+[mark=none,domain=0:14,samples=200]{1-exp(-0.3*x)/sqrt(1-0.09)*sin(deg(sqrt(1-0.09)*x)+deg(acos(0.3)))}; \addlegendentry{$\zeta=0.3$}
\addplot+[mark=none,domain=0:14,samples=150]{1-exp(-0.7071*x)/sqrt(0.5)*sin(deg(sqrt(0.5)*x)+45)}; \addlegendentry{$\zeta=0.707$}
\addplot+[mark=none,domain=0:14,samples=100]{1-(1+x)*exp(-x)}; \addlegendentry{$\zeta=1$}
\addplot+[mark=none,domain=0:14,samples=100]{1+(-3.7321*exp(-0.26795*x)+0.26795*exp(-3.7321*x))/(3.7321-0.26795)}; \addlegendentry{$\zeta=2$}
\end{axis}
\end{tikzpicture}
\end{column}
\begin{column}{.38\textwidth}
\centering
\begin{tikzpicture}[scale=.9]
\draw[-{Stealth},mkGray] (-2.6,0) -- (0.5,0) node[right,font=\tiny]{Re};
\draw[-{Stealth},mkGray] (0,-1.5) -- (0,1.5) node[above,font=\tiny]{Im};
\draw[mkSky,dotted,line width=.6pt] (0,0) circle (1.2);
\foreach \z/\c in {0.1/mkOrange,0.3/mkAmber,0.707/mkBlue}{
 \pgfmathsetmacro{\re}{-1.2*\z} \pgfmathsetmacro{\im}{1.2*sqrt(1-\z*\z)}
 \node[text=\c,font=\small] at (\re,\im) {$\times$}; \node[text=\c,font=\small] at (\re,-\im) {$\times$};}
\node[text=mkMid,font=\small] at (-1.2,0) {$\times$};
\node[text=mkSky,font=\small] at (-0.32,0) {$\times$}; \node[text=mkSky,font=\small] at (-2.4,0.0) {$\times$};
\draw[mkGray,dashed] (0,0) -- (135:1.35);
\node[font=\tiny,text=mkGray] at (-1.3,1.25) {$\zeta=0.707$};
\node[font=\tiny,text=mkGray] at (0.75,-1.1) {$|s|=\omega_n$};
\end{tikzpicture}

\footnotesize
\begin{itemize}\setlength\itemsep{1pt}
\item Poles $s=-\zeta\omega_n\pm j\omega_n\sqrt{1-\zeta^2}$.
\item Overshoot $M_p=e^{-\pi\zeta/\sqrt{1-\zeta^2}}$: 73\,\%, 37\,\%, 4.3\,\%.
\item Lightly damped mechanics ($\zeta<0.1$) is the norm in servo structures.
\end{itemize}
\end{column}
\end{columns}
\end{frame}

% ---------------------------------------------------------------------
\begin{frame}{Rotational mechanical systems}
\begin{columns}[T]
\begin{column}{.48\textwidth}
\centering
\begin{tikzpicture}[scale=.9]
\fill[pattern=north east lines,pattern color=mkGray] (-0.3,0.2) rectangle (0,2.2); \draw[mkGray,line width=1pt] (0,0.2) -- (0,2.2);
\draw[mkNavy,line width=2.2pt] (0,1.2) -- (2.4,1.2);
\node[font=\scriptsize] at (1.2,1.55) {shaft $k_\theta$};
\fill[mkSky!70,draw=mkBlue] (2.4,0) rectangle (2.9,2.4); \node[font=\scriptsize] at (2.65,2.65) {$J$};
\draw[mkGray,line width=1.5pt] (2.9,1.2) -- (3.6,1.2);
\fill[mkGray!40] (3.6,0.8) rectangle (4.0,1.6); \node[font=\tiny] at (3.8,0.55) {$B$ (bearing)};
\draw[-{Stealth},mkOrange,line width=1pt] (3.2,1.8) arc (100:260:0.25 and 0.6);
\node[font=\scriptsize,text=mkOrange] at (3.35,2.15) {$T(t)$};
\draw[-{Stealth},mkNavy] (2.65,-0.2) arc (-90:30:0.35) ;
\node[font=\scriptsize] at (3.2,-0.3) {$\theta$};
\end{tikzpicture}
\begin{eqbox}\small
$J\,\ddot\theta + B\,\dot\theta + k_\theta\,\theta = T(t)$\\[1mm]
$\dfrac{\Theta(s)}{T(s)}=\dfrac{1}{Js^2+Bs+k_\theta}$
\end{eqbox}
\end{column}
\begin{column}{.50\textwidth}
\footnotesize
\begin{itemize}\setlength\itemsep{2pt}
\item Euler's law: $\sum T = J\ddot\theta$ about a fixed axis.
\item Torsional stiffness of a solid shaft: $k_\theta=\dfrac{G\,I_p}{l}=\dfrac{G\,\pi d^4}{32\,l}$.
\item Example: steel ($G=\SI{80}{\giga\pascal}$), $d=\SI{10}{\milli\metre}$, $l=\SI{100}{\milli\metre}$\\
$\Rightarrow k_\theta=\SI{785}{\newton\metre\per\radian}$.
\item Disc inertia $J=\tfrac12 m r^2$, hollow cylinder $J=\tfrac12 m(r_o^2+r_i^2)$.
\item Couplings, gear meshes and belts are springs too: series stiffness $1/k=\sum 1/k_i$ (softest element dominates).
\item Servo motor without load: $k_\theta=0$ $\Rightarrow$ $\Theta/T=1/(s(Js+B))$ -- a double integrator for $B\to0$.
\end{itemize}
\end{column}
\end{columns}
\end{frame}

% ---------------------------------------------------------------------
\begin{frame}{Electrical RLC network}
\begin{columns}[T]
\begin{column}{.46\textwidth}
\centering
\begin{circuitikz}[scale=.85,transform shape,font=\small]
\draw (0,0) to[V,v<=$v(t)$,color=mkNavy] (0,3)
 to[R,l=$R$,i>^=$i$,color=mkNavy] (3,3)
 to[L,l=$L$,color=mkNavy] (5,3)
 to[C,l=$C$,v^<=$v_C$,color=mkOrange] (5,0) -- (0,0);
\end{circuitikz}
\begin{eqbox}
$L\,\ddot q + R\,\dot q + \dfrac{1}{C}\,q = v(t)$, \quad $i=\dot q$\\[1mm]
$\dfrac{V_C(s)}{V(s)}=\dfrac{1}{LCs^2+RCs+1}$
\end{eqbox}
\end{column}
\begin{column}{.52\textwidth}
\footnotesize
\begin{itemize}\setlength\itemsep{2pt}
\item KVL around the loop: $v=v_R+v_L+v_C$; constitutive laws $v_R=Ri$, $v_L=L\,di/dt$, $i=C\,dv_C/dt$.
\item Same structure as the mass--spring--damper:\\
$\omega_n=\dfrac{1}{\sqrt{LC}}$, \quad $\zeta=\dfrac{R}{2}\sqrt{\dfrac{C}{L}}$.
\item In drives: DC-link $LC$ input filter, sine/d$v$/d$t$ output filters, cable resonances.
\item Example: $L=\SI{1}{\milli\henry}$, $C=\SI{470}{\micro\farad}$, $R=\SI{0.1}{\ohm}$\\
$\Rightarrow f_n=\SI{232}{\hertz}$, $\zeta=0.034$ -- a lightly damped filter.
\item Motor winding alone ($C\to\infty$): first-order $1/(Ls+R)$, $\tau_e=L/R$.
\end{itemize}
\end{column}
\end{columns}
\end{frame}

% ---------------------------------------------------------------------
\begin{frame}{Analogies and energy-based modelling}
\begin{columns}[T]
\begin{column}{.56\textwidth}
\scriptsize
\renewcommand{\arraystretch}{1.15}
\begin{tabular}{@{}llll@{}}
\toprule
\textbf{Translational} & \textbf{Rotational} & \textbf{Force--voltage} & \textbf{Force--current}\\
\midrule
force $F$ & torque $T$ & voltage $v$ & current $i$\\
velocity $\dot x$ & speed $\omega$ & current $i$ & voltage $v$\\
mass $m$ & inertia $J$ & inductance $L$ & capacitance $C$\\
damper $b$ & damper $B$ & resistance $R$ & conductance $1/R$\\
spring $k$ & torsion $k_\theta$ & elastance $1/C$ & reciprocal ind. $1/L$\\
$\tfrac12 m\dot x^2$ & $\tfrac12 J\omega^2$ & $\tfrac12 L i^2$ & $\tfrac12 C v^2$\\
$\tfrac12 k x^2$ & $\tfrac12 k_\theta\theta^2$ & $\tfrac12 q^2/C$ & $\tfrac12 \lambda^2/L$\\
\bottomrule
\end{tabular}\\[1mm]
{\scriptsize Product of \textbf{effort} and \textbf{flow} is always \hl{power}: $Fv$, $T\omega$, $vi$.}
\end{column}
\begin{column}{.42\textwidth}
\begin{defbox}[Lagrange's equations]
\small
$\mathcal{L}=T_\text{kin}-V_\text{pot}$, \; $D$ = Rayleigh dissipation
\[
\frac{d}{dt}\frac{\partial\mathcal{L}}{\partial\dot q_j}-\frac{\partial\mathcal{L}}{\partial q_j}+\frac{\partial D}{\partial\dot q_j}=Q_j
\]
\end{defbox}
\footnotesize
MSD: $\mathcal{L}=\tfrac12m\dot x^2-\tfrac12kx^2$, $D=\tfrac12b\dot x^2$, $Q=F$\\
$\Rightarrow m\ddot x+b\dot x+kx=F$ \;\textcolor{mkGreen}{\checkmark}
\begin{itemize}\setlength\itemsep{1pt}
\item Energy view gives the model \hl{without free-body diagrams} -- useful for multi-body and coupled systems.
\item Same Lagrangian approach with $q$ = charge yields the RLC equation.
\end{itemize}
\end{column}
\end{columns}
\end{frame}

% ---------------------------------------------------------------------
\begin{frame}{Electromechanical coupling: the motor as a gyrator}
\begin{columns}[T]
\begin{column}{.52\textwidth}
\centering
\resizebox{\linewidth}{!}{\begin{tikzpicture}[font=\scriptsize]
\node[mkblk,minimum width=17mm,minimum height=13mm] (e) {electrical\\ $R,\;L$};
\node[mkblk,minimum width=17mm,minimum height=13mm,fill=mkCream,draw=mkOrange,right=9mm of e] (g) {gyrator\\ $k=k_t=k_e$};
\node[mkblk,minimum width=17mm,minimum height=13mm,right=9mm of g] (m) {mechanical\\ $J,\;B$};
\draw[mksig,{Stealth}-{Stealth}] (e) -- node[above]{$e,\,i$} (g);
\draw[mksig,{Stealth}-{Stealth}] (g) -- node[above]{$T,\,\omega$} (m);
\draw[mksig] ($(e.west)+(-7mm,0)$) -- node[above]{$v$} (e);
\draw[mksig] ($(m.east)+(7mm,0)$) -- node[above]{$T_L$} (m);
\node[below=2mm of g,align=center,text=mkNavy] {$T=k\,i,\quad e=k\,\omega$\\ (effort $\leftrightarrow$ flow swapped)};
\end{tikzpicture}}

\vspace{2mm}
\begin{eqbox}\scriptsize
State $\bm{x}=[i\;\;\omega]^\mathsf{T}$:\\
$\dfrac{d}{dt}\begin{bmatrix} i\\ \omega\end{bmatrix}=
\begin{bmatrix}-R/L & -k/L\\ k/J & -B/J\end{bmatrix}\begin{bmatrix} i\\ \omega\end{bmatrix}
+\begin{bmatrix}1/L & 0\\ 0 & -1/J\end{bmatrix}\begin{bmatrix} v\\ T_L\end{bmatrix}$
\end{eqbox}
\end{column}
\begin{column}{.46\textwidth}
\footnotesize
\textbf{Energy balance} of the coupled system:
\[
v\,i=\underbrace{R i^2 + B\omega^2}_{\text{dissipated}}+\frac{d}{dt}\underbrace{\Big(\tfrac12Li^2+\tfrac12J\omega^2\Big)}_{\text{stored}}+T_L\,\omega
\]
\begin{itemize}\setlength\itemsep{1pt}
\item The coupling terms ($\pm k$) form a \hl{skew-symmetric} (power-conserving) interconnection: they exchange energy, never create or dissipate it.
\item Seen from the terminals, inertia $J$ looks like a \hl{capacitance} $J/k^2$ and damping $B$ like a conductance $B/k^2$.
\item Two energy stores $\Rightarrow$ 2nd-order system, as found in Session 1.3.
\item State-space form in detail: Day 2 (session 2.2).
\end{itemize}
\end{column}
\end{columns}
\end{frame}

% ---------------------------------------------------------------------
\begin{frame}{Preview: two-inertia system (motor -- compliant shaft -- load)}
\begin{columns}[T]
\begin{column}{.42\textwidth}
\centering
\begin{tikzpicture}[scale=.85]
\fill[mkBlue] (0,0) rectangle (1,1.4); \node[white,font=\scriptsize] at (0.5,0.7) {$J_m$};
\draw[-{Stealth},mkOrange,line width=1pt] (-0.4,1.2) arc (110:250:0.25 and 0.5); \node[font=\scriptsize,text=mkOrange] at (-0.5,1.5) {$T_m$};
\draw[decorate,decoration={coil,aspect=0.5,segment length=2mm,amplitude=2mm},mkNavy] (1,0.95) -- (2.8,0.95);
\node[font=\scriptsize] at (1.9,1.4) {$k$};
\draw[mkNavy] (1,0.45) -- (1.6,0.45); \draw[mkNavy] (1.6,0.3) -- (1.6,0.6); \draw[mkNavy] (1.5,0.25)--(2.0,0.25)--(2.0,0.65)--(1.5,0.65); \draw[mkNavy] (2.0,0.45) -- (2.8,0.45);
\node[font=\scriptsize] at (1.8,0.0) {$c$};
\fill[mkOrange] (2.8,0) rectangle (3.8,1.4); \node[white,font=\scriptsize] at (3.3,0.7) {$J_L$};
\node[font=\scriptsize] at (0.5,-0.3) {$\theta_m$}; \node[font=\scriptsize] at (3.3,-0.3) {$\theta_L$};
\end{tikzpicture}
\begin{eqbox}\footnotesize
$J_m\ddot\theta_m=T_m-k(\theta_m-\theta_L)-c(\dot\theta_m-\dot\theta_L)$\\
$J_L\ddot\theta_L=k(\theta_m-\theta_L)+c(\dot\theta_m-\dot\theta_L)-T_L$
\end{eqbox}
\footnotesize
$\dfrac{\Omega_m}{T_m}=\dfrac{J_Ls^2+cs+k}{s\,[J_mJ_Ls^2+c(J_m+J_L)s+k(J_m+J_L)]}$
\end{column}
\begin{column}{.56\textwidth}
\begin{tikzpicture}
\begin{axis}[mkplot,width=7.8cm,height=4.6cm,xmode=log,xmin=10,xmax=2000,ymin=-30,ymax=45,
  xlabel={frequency (\si{\hertz})},ylabel={$|\Omega_m/T_m|$ (dB)},legend pos=north east]
\addplot[mkBlue,thick,domain=1:3.301,samples=500] ({10^x},
 {20*log10( sqrt((300-6e-4*(2*pi*10^x)^2)^2+(0.02*2*pi*10^x)^2) / ( (2*pi*10^x)*sqrt((0.24-1.2e-7*(2*pi*10^x)^2)^2+(1.6e-5*2*pi*10^x)^2) ) )});
\addlegendentry{two-inertia}
\addplot[mkOrange,thick,dashed,domain=1:3.301,samples=50] ({10^x},{20*log10(1/(8e-4*2*pi*10^x))});
\addlegendentry{rigid $1/((J_m+J_L)s)$}
\addplot[mkSky,thick,dotted,domain=1:3.301,samples=50] ({10^x},{20*log10(1/(2e-4*2*pi*10^x))});
\addlegendentry{motor only $1/(J_ms)$}
\node[font=\tiny,text=mkNavy] at (axis cs:112,-26) {anti-res.\ 113\,Hz};
\node[font=\tiny,text=mkNavy] at (axis cs:330,30) {res.\ 225\,Hz};
\end{axis}
\end{tikzpicture}\\
{\tiny $J_m=\SI{2e-4}{}$, $J_L=\SI{6e-4}{\kilo\gram\metre\squared}$, $k=\SI{300}{\newton\metre\per\radian}$, $c=\SI{0.02}{\newton\metre\second\per\radian}$}
\footnotesize
\begin{itemize}\setlength\itemsep{1pt}
\item Anti-resonance $\omega_{ar}=\sqrt{k/J_L}$; resonance $\omega_r=\sqrt{k(J_m+J_L)/(J_mJ_L)}$.
\item Limits achievable servo bandwidth; motivates notch filters (Day 3, Day 5).
\end{itemize}
\end{column}
\end{columns}
\end{frame}

% ---------------------------------------------------------------------
\begin{frame}{Session 2.1 (part 1) -- Summary / key take-aways}
\begin{columns}[T]
\begin{column}{.54\textwidth}
\begin{itemize}
\item Lumped models: stores of kinetic and potential energy plus dissipators; each store adds a state.
\item Mass--spring--damper, torsional system and RLC circuit share one 2nd-order form with $\omega_n,\zeta$.
\item Analogies (force--voltage, force--current) let us reuse intuition and tools across domains.
\item The motor is a gyrator: $T=ki$, $e=k\omega$; it couples the RL circuit to the inertia.
\item Compliance between motor and load creates resonance and anti-resonance.
\end{itemize}
\end{column}
\begin{column}{.42\textwidth}
\begin{keybox}[Remember]
\small
$m\ddot x+b\dot x+kx=F$\\[1mm]
$\omega_n=\sqrt{k/m},\;\;\zeta=b/(2\sqrt{km})$\\[1mm]
$\omega_{ar}=\sqrt{k/J_L}$\\[1mm]
$\omega_r=\sqrt{k(J_m+J_L)/(J_mJ_L)}$
\end{keybox}
\end{column}
\end{columns}
\end{frame}

% =====================================================================
\section*{Wrap-up}
% =====================================================================
\begin{frame}{Day 1 summary}
\begin{columns}[T]
\begin{column}{.5\textwidth}
\footnotesize
\begin{defbox}[Machines]
Lorentz/Faraday $\Rightarrow$ $k_t=k_e$; reluctance $\tfrac12 i^2dL/d\theta$.\\
DC, BLDC (six-step), PMSM (sinusoidal), induction (slip), stepper (open loop).
\end{defbox}
\begin{defbox}[Models]
$v=Ri+L\dot i+k_e\omega$, \; $J\dot\omega=k_ti-B\omega-T_L$\\
$\tau_e=L/R$, $\tau_m=RJ/(k_tk_e)$; friction is nonlinear.
\end{defbox}
\end{column}
\begin{column}{.48\textwidth}
\footnotesize
\begin{defbox}[Transmissions and sizing]
$J/N^2$, $m(p/2\pi)^2$, $mr^2$; $N_\text{opt}=\sqrt{J_L/J_m}$.\\
Peak in intermittent zone, $T_\text{rms}$ in continuous zone; match $V_\text{DC}$, $k_e$, currents.
\end{defbox}
\begin{defbox}[Dynamic modelling]
Lumped elements, FBD, analogies, Lagrange; 2nd-order prototype $\omega_n,\zeta$; two-inertia resonance.
\end{defbox}
\end{column}
\end{columns}
\vspace{2mm}
\begin{keybox}[One sentence]
A servo axis is an \hl{energy converter} (motor) plus a \hl{mechanical network} (transmission, load) -- both described by the same lumped, energy-based equations.
\end{keybox}
\end{frame}

% ---------------------------------------------------------------------
\begin{frame}{Preview of Day 2: System dynamics and control}
\begin{columns}[T]
\begin{column}{.55\textwidth}
\small
\renewcommand{\arraystretch}{1.15}
\begin{tabular}{@{}lp{5.2cm}r@{}}
\toprule
\textbf{Code} & \textbf{Session} & \textbf{Time}\\
\midrule
2.1 & Modelling (cont.): nonlinear effects and linearisation & 0.5\,h\\
2.2 & Transfer functions and state space & 1.0\,h\\
2.3 & System identification and frequency response & 1.0\,h\\
2.4 & PID and basic control structures & 1.5\,h\\
2.5 & Stability and performance analysis & 1.0\,h\\
\bottomrule
\end{tabular}
\end{column}
\begin{column}{.42\textwidth}
\begin{notebox}
\textbf{Bring along from today:}\\
the DC-motor model and its transfer function, the two time constants, the friction nonlinearity, and the 2nd-order prototype with $\omega_n$ and $\zeta$ -- they are the plants we will linearise, identify and control tomorrow.
\end{notebox}
\end{column}
\end{columns}
\end{frame}

% ---------------------------------------------------------------------
\begin{frame}{Figure credits}
\scriptsize
\renewcommand{\arraystretch}{1.1}
\begin{tabularx}{\textwidth}{@{}l l X l@{}}
\toprule
\textbf{Slide} & \textbf{File (Wikimedia Commons)} & \textbf{Author} & \textbf{License}\\
\midrule
Brushed DC motor & Motor disassembeled.JPG & Jfer91 & CC BY-SA 3.0\\
BLDC / PMSM & Floppy drive spindle motor open.jpg & Sebastian Koppehel & CC BY 3.0\\
Induction motor & Wirnik by Zureks.jpg & Zureks & CC BY-SA 3.0\\
Stepper motor & Positioner Stepper Motor parts.jpg & tony\_duell & CC BY 2.0\\
Linear actuators & Linear motor by Zureks.jpg & Zureks & CC BY-SA 4.0\\
Ball screw & BallScrewNut-showing-balls-recirculating-tubes.jpg & Glenn McKechnie (Graibeard) & CC BY-SA 3.0\\
Gearboxes & Epicyclic gear ratios.png & Wapcaplet & CC BY-SA 3.0\\
Gearboxes & Harmonic Drive AG strain wave gear set.jpg & Pieceofmetalwork & CC BY-SA 4.0\\
\bottomrule
\end{tabularx}
{\tiny Source pages: \texttt{https://commons.wikimedia.org/wiki/File:<file name>} (see \texttt{figures/CREDITS.tsv}).}
\vspace{2mm}
All other drawings, plots and diagrams: own work for this course (TikZ / pgfplots / circuitikz).
\end{frame}

% ---------------------------------------------------------------------
\begin{frame}{References}
\scriptsize
\begin{itemize}\setlength\itemsep{1pt}
\item A.\,E. Fitzgerald, C. Kingsley, S.\,D. Umans, \textit{Electric Machinery}, 6th ed., McGraw-Hill, 2003.
\item R. Krishnan, \textit{Electric Motor Drives: Modeling, Analysis, and Control}, Prentice Hall, 2001.
\item R. Krishnan, \textit{Permanent Magnet Synchronous and Brushless DC Motor Drives}, CRC Press, 2010.
\item N. Mohan, \textit{Electric Drives: An Integrative Approach}, MNPERE, 2003.
\item A. Hughes, B. Drury, \textit{Electric Motors and Drives}, 4th ed., Newnes, 2013.
\item W. Leonhard, \textit{Control of Electrical Drives}, 3rd ed., Springer, 2001.
\item J.\,R. Hendershot, T.\,J.\,E. Miller, \textit{Design of Brushless Permanent-Magnet Machines}, Motor Design Books, 2010.
\item T. Kenjo, A. Sugawara, \textit{Stepping Motors and Their Microprocessor Controls}, 2nd ed., Oxford Univ. Press, 1994.
\item H. Gro\ss, J. Hamann, G. Wiegärtner, \textit{Electrical Feed Drives in Automation}, Publicis MCD (Siemens), 2001.
\item G. Ellis, \textit{Control System Design Guide}, 4th ed., Butterworth-Heinemann, 2012.
\item K. Ogata, \textit{Modern Control Engineering}, 5th ed., Prentice Hall, 2010.
\item G.\,F. Franklin, J.\,D. Powell, A. Emami-Naeini, \textit{Feedback Control of Dynamic Systems}, 7th ed., Pearson, 2015.
\item D.\,C. Karnopp, D.\,L. Margolis, R.\,C. Rosenberg, \textit{System Dynamics: Modeling, Simulation, and Control of Mechatronic Systems}, 5th ed., Wiley, 2012.
\item B. Armstrong-H\'elouvry, P. Dupont, C. Canudas de Wit, ``A survey of models, analysis tools and compensation methods for the control of machines with friction'', \textit{Automatica} 30(7), 1994.
\item IEC 60034-1, \textit{Rotating electrical machines -- Part 1: Rating and performance}; IEC 61800 series, \textit{Adjustable speed electrical power drive systems}.
\end{itemize}
\end{frame}
\end{document}
